% !TeX root = ../../Book1.tex
% 纲要标题：## 附录 F　Gröbner 基初步
% 纲要位置：工作记录/计划纲要.md，第 951 行。
% 纲要细目（写作范围）：
% > 第10—12章之后可插读。承担多元除法、Gröbner基存在性、正规式、Buchberger判据与算法、约化基唯一性、基本消去和标准单项式基的主要初等证明；有限性精确回引第12.4节。第二卷非交换专题建议先读本附录F.1—F.3，第三卷第8章为独立阅读重新完整建立基础理论，再发展进阶应用。



\chapter{Gröbner 基初步}
\label{b1:app:F}

\begin{chapterabstract}
多元多项式除法一般没有唯一余式，理想的任意生成组也未必适合判断一个多项式是否属于理想。本附录从这两个困难出发，建立全局单项式序、Gröbner 基及 Buchberger 算法，证明正规式、约化基与基本消去定理。最后用标准单项式给商环选择坐标，计算乘法、单位与零因子。所有核心论证都在域上的多项式环内完成，并通过可直接展开核验的等式保存计算依据。
\end{chapterabstract}

\begin{learninggoals}
\begin{itemize}
\item 区分首项、单项式序与全次数，执行多元除法并说明余式为何可能依赖选择。
\item 从首项理想理解 Gröbner 基，证明存在性、正规式唯一性和成员判定。
\item 解释 S 多项式的抵消机制，证明并执行 Buchberger 算法，求出约化基。
\item 选取适当的字典序消去变量，区分消去关系、代回及解所在的域。
\item 用标准单项式计算商环，检验单位并辨认点集无法直接反映的幂零元素。
\end{itemize}
\end{learninggoals}

读完第10—12章即可开始本附录。所需知识是理想、商环、一元除法、多元多项式，以及\cref{b1:thm:hilbert-basis}和\cref{b1:prop:noetherian-acc}；商环的基沿用线性代数中的含义，使用时会写出唯一的线性组合表示。一般定理适用于任意域，算法实际执行则要求系数域的四则运算和零判定能够精确完成。主要算例取 \(\mathbb Q\) 与素域，并明确指出改变系数域的影响。

\ProseParagraphBreak
本附录是第二卷第20.6节非交换 Gröbner 专题之前的交换情形入门；第三卷第8章将重述必要的定义与结果，再发展理想运算、几何消去和零维计算。初次阅读可先顺读前三节，亲自完成一个从生成组到约化基的计算，再用后两节解释输出的代数意义。

\ProseParagraphBreak
进一步阅读可参阅 Cox、Little 与 O'Shea 的《Ideals, Varieties, and Algorithms》\cite[Chapter 2, \S\S 2--7; Chapter 3, \S 1; Chapter 5, \S\S 2--3]{CoxLittleOShea2015IdealsVarietiesAlgorithms}。这些位置分别讨论单项式序与 Gröbner 基、消去，以及商环中的计算；本附录将基本证明完整写出，具体计算不以特定软件为前提。

% !TeX root = ../../Book1.tex
% 纲要标题：### F.1 单项式序与多元除法
% 纲要位置：工作记录/计划纲要.md，第 955 行。

% 纲要要点（供目录核对）：
% #### F.1.1 单项式的整除与次序
% * 多重指数与单项式整除，单项式的最大公因式和最小公倍式。
% * 全局单项式序的全序、良序和乘法相容条件。
% * 以字典序、分次字典序为主要例子，比较它们选择首项的方式；分次逆字典序只作简短补充。
% #### F.1.2 首项与多元除法
% * 首单项式、首项系数和首项的区别；固定记号与零多项式的处理约定。
% * 写出除以一组多项式的逐步算法，保存等式 \(f=\sum_iq_ig_i+r\)。
% * 用首单项式的严格下降证明终止，说明余式中的单项式均不能被任何除式的首单项式整除。
% #### F.1.3 余式为何依赖除式顺序
% * 在 \(\mathbb Q[x,y]\) 中固定字典序 \(x>y\)，用 \(xy-1,y^2-1\) 除 \(xy^2\)，分别得到余式 \(y\) 与 \(x\)。
% * 解释零余式总能提供理想成员证据，而对任意生成组得到非零余式，不能据此断定不属于理想。
% * 引出问题：怎样选取生成组，才能让余式唯一并准确判断理想成员？

\section{单项式序与多元除法}
\label{b1:sec:F01}

一元多项式除法每次消去次数最高的一项，次数便严格下降。到了两个变量，\(x^2\) 与 \(y^3\) 哪一项应当先消去，却不能仅由“多项式”这个概念决定。必须先为单项式选定次序，再问当前首项能否被某个除式的首项整除。这样的除法可以完成，但它最初并不具有一元除法的全部优点：同一个被除式，换一个除式顺序，就可能得到不同的余式。

\ProseParagraphBreak
本附录始终令 \(k\) 为域、\(n\geq1\) 为整数，并置 \(R=k[x_1,\ldots,x_n]\)。一般结论不限制 \(k\) 的特征；实际运行算法时，还要求能够在 \(k\) 中进行精确四则运算并判定一个系数是否为零。有理数域和给定的有限域满足这些要求。

\subsection{单项式的整除与次序}
\label{b1:subsec:F01:01}

沿用\cref{b1:def:multivariate-degree}前的记号，若 \(\alpha=(\alpha_1,\ldots,\alpha_n)\in\mathbb N^n\)，则
\[
x^\alpha=x_1^{\alpha_1}\cdots x_n^{\alpha_n},\qquad
|\alpha|=\alpha_1+\cdots+\alpha_n.
\]
这里单项式的系数取为 \(1\)；\(c x^\alpha\)，其中 \(c\in k\) 非零，称为多项式的一项。单项式的整除可以逐个变量检查：
\[
x^\alpha\mid x^\beta
\quad\Longleftrightarrow\quad
\alpha_i\leq\beta_i\quad(1\leq i\leq n).
\]
在满足这些不等式时，商是 \(x^{\beta-\alpha}\)。因此最大公因式和最小公倍式分别由各分量的最小值和最大值给出：
\[
\begin{aligned}
\gcd(x^\alpha,x^\beta)
 &=\prod_{i=1}^n x_i^{\min(\alpha_i,\beta_i)},\\
\operatorname{lcm}(x^\alpha,x^\beta)
 &=\prod_{i=1}^n x_i^{\max(\alpha_i,\beta_i)}.
\end{aligned}
\]
例如 \(x^2y\) 与 \(xy^3\) 的最大公因式是 \(xy\)，最小公倍式是 \(x^2y^3\)。整除只是偏序：\(x\) 与 \(y\) 互不整除；消项算法需要另外选一个能比较任意两个单项式的次序。

\begin{definition}\label{b1:def:F-monomial-order}
设 \(k\) 为域，\(n\geq1\) 为整数，\(R=k[x_1,\ldots,x_n]\)。若 \(R\) 的单项式集合上的全序 \(\preceq\) 满足以下条件：
\begin{enumerate}
\item 任一非空单项式集合都有最小元，即这个全序是良序；
\item 对任意单项式 \(u,v,w\)，若 \(u\prec v\)，则 \(uw\prec vw\)。
\end{enumerate}
则称 \(\preceq\) 为 \(R\) 上的\term[danxiangshi-xu]{单项式序}[monomial order]。本附录只采用这样的全局单项式序；符号 \(u\prec v\) 表示 \(u\preceq v\) 且 \(u\ne v\)。
\end{definition}

这些条件蕴含 \(1\) 是最小单项式。否则若 \(u\prec1\)，乘法相容性就给出无穷严格下降链 \(1\succ u\succ u^2\succ\cdots\)，其各项组成的集合没有最小元，与良序矛盾。因而只要 \(u\mid v\)，就有 \(u\preceq v\)；若 \(u\ne v\)，这个不等式严格成立。反过来，两个单项式在所选次序下可比，并不意味着它们存在整除关系。

\ProseParagraphBreak
最直接的例子是\term[zidian-xu]{字典序}[lexicographic order]。规定变量优先级为 \(x_1\succ\cdots\succ x_n\)，比较 \(x^\alpha\) 和 \(x^\beta\) 时，从第一个坐标开始找出 \(\alpha_j\ne\beta_j\) 的位置，并规定
\[
x^\alpha\prec_{\mathrm{lex}}x^\beta
\quad\Longleftrightarrow\quad \alpha_j<\beta_j.
\]
这与按首个不同字母排列单词相似，但这里比较的是各变量的指数。它确实是良序：在任一非空指数集合中，先把第一坐标限制为其中的最小值，再在剩下的集合中把第二坐标限制为最小值，经过有限个坐标后得到字典序最小的指数。给两组指数同时加上 \(\gamma\)，首个不同位置及其大小关系不变，所以乘法相容性也成立。字典序不按全次数排列；在 \(k[x,y]\) 中取 \(x\succ y\) 时，对任意 \(m\geq1\) 都有 \(x\succ y^m\)。

\ProseParagraphBreak
若希望先比较全次数，可以采用\term[fenci-zidian-xu]{分次字典序}[graded lexicographic order]：先比较 \(|\alpha|\) 和 \(|\beta|\)，全次数相同时再按字典序比较。任何非空单项式集合先有最小全次数，而固定全次数的单项式只有有限多个，因此这是良序。乘以同一个单项式时，两边全次数同时增加同一个数；若原来全次数相等，字典序的比较又保持不变。所以它也是单项式序。例如 \(x+y^2\) 在字典序 \(x\succ y\) 下以 \(x\) 为最大单项式，在分次字典序下却以 \(y^2\) 为最大单项式。

\ProseParagraphBreak
计算中还常用\term[fenci-ni-zidian-xu]{分次逆字典序}[graded reverse lexicographic order]：先比较全次数；全次数相同时，从最后一个坐标向前寻找首个不同位置 \(j\)，并规定指数较小者的单项式较大，即 \(\alpha_j<\beta_j\) 时 \(x^\alpha\succ x^\beta\)。在每一固定全次数上，这就是对反向排列并改变符号的指数作字典序比较，所以是全序。对非空单项式集合，先选其中最小的全次数，再在该次数的有限集合中取最小元，便得到良序性；加上同一个指数向量仍保持比较，故乘法相容。例如在 \(k[x,y,z]\) 中，\(x^2z\) 与 \(xy^2\) 的全次数都为三：字典序和分次字典序把前者排得更大，分次逆字典序则把后者排得更大。后文证明只使用\cref{b1:def:F-monomial-order}的两项条件；每个算例则会明确所取次序。

\subsection{首项与多元除法}
\label{b1:subsec:F01:02}

固定一个单项式序。非零多项式只含有限多个单项式，因而其中有唯一的最大者。把它与所带系数分开记，可以避免在整除和消项时混淆两类数据。

\begin{definition}\label{b1:def:F-leading-data}
设 \(k\) 为域，\(n\geq1\) 为整数，\(R=k[x_1,\ldots,x_n]\) 配备一个单项式序。对非零多项式 \(f=\sum_\alpha c_\alpha x^\alpha\in R\)，设 \(x^\beta\) 是满足 \(c_\beta\ne0\) 的单项式中最大者。令
\[
\operatorname{LM}(f)=x^\beta,\qquad
\operatorname{LC}(f)=c_\beta,\qquad
\operatorname{LT}(f)=c_\beta x^\beta
\]
分别称 \(\operatorname{LM}(f)\)、\(\operatorname{LC}(f)\) 与 \(\operatorname{LT}(f)\) 为 \(f\) 的\term[shou-danxiangshi]{首单项式}[leading monomial]、首项系数与\term[shouxiang]{首项}[leading term]。这些记号总是相对于当前固定的单项式序。零多项式不定义首单项式和首项系数；含有这些记号的比较均只针对非零多项式。
\end{definition}

\termidx[shouxiang-xishu]{首项系数}
\symidx{\(\operatorname{LM}(f)\)：多项式的首单项式}
\symidx{\(\operatorname{LC}(f)\)：多项式的首项系数}
\symidx{\(\operatorname{LT}(f)\)：多项式的首项}

例如在 \(\mathbb Q[x,y]\) 中，取字典序 \(x\succ y\)，则 \(f=3x^2y-5xy^4+2\) 满足 \(\operatorname{LM}(f)=x^2y\)、\(\operatorname{LC}(f)=3\)、\(\operatorname{LT}(f)=3x^2y\)；若换为分次字典序，三者依次变成 \(xy^4\)、\(-5\)、\(-5xy^4\)。因此“首”表示所选单项式序中的最大位置，不总是全次数最高的位置。

\begin{proposition}\label{b1:prop:F-leading-product}
设 \(k\) 为域，\(n\geq1\) 为整数，\(R=k[x_1,\ldots,x_n]\) 配备一个单项式序。若 \(f,g\in R\) 均非零，则
\[
\operatorname{LM}(fg)=\operatorname{LM}(f)\operatorname{LM}(g),\qquad
\operatorname{LT}(fg)=\operatorname{LT}(f)\operatorname{LT}(g).
\]
若 \(f+g\ne0\)，则
\[
\operatorname{LM}(f+g)\preceq
\max\Bigl\{\operatorname{LM}(f),\operatorname{LM}(g)\Bigr\}.
\]
当 \(f\)、\(g\) 的首单项式不同，和式中的不等式取等号。
\end{proposition}
\begin{proof}
记 \(u=\operatorname{LM}(f)\)、\(v=\operatorname{LM}(g)\)。若 \(u'\)、\(v'\) 分别来自 \(f\)、\(g\)，则 \(u'\preceq u\)、\(v'\preceq v\)，所以 \(u'v'\preceq uv\)。除 \((u',v')=(u,v)\) 外，至少一个比较严格；乘法相容性使乘积比较也严格。因此 \(uv\) 只来自两个首项的乘积，其系数是两个非零首项系数的积，在域中不为零。这就证明前两式。相加不会引入原来两式都没有的单项式；若两个首单项式不同，较大者也没有另一项可以与它抵消，故得到最后的断言。
\end{proof}

现在给定有序的一组非零多项式 \(g_1,\ldots,g_s\)。若当前多项式 \(p\) 的首单项式被 \(\operatorname{LM}(g_i)\) 整除，就可以减去
\[
\frac{\operatorname{LT}(p)}{\operatorname{LT}(g_i)}g_i
=\frac{\operatorname{LC}(p)}{\operatorname{LC}(g_i)}
 \frac{\operatorname{LM}(p)}{\operatorname{LM}(g_i)}g_i
\]
来消去首项。右边两个分式分别是域中的系数商和整除成立时的单项式商；并没有在 \(R\) 中对任意多项式作除法。

\ProseParagraphBreak
把这一步反复应用，就得到下面的\term[duoyuan-chufa]{多元除法}[multivariate division]。若当前首项不能消去，算法并不立刻停下，而是把这一项移入余式，再处理剩余部分。

\begin{theorem}[多元除法]\label{b1:thm:F-division}
设 \(k\) 为域，\(n\geq1\) 为整数，\(R=k[x_1,\ldots,x_n]\) 配备一个单项式序。对 \(f\in R\) 和有限个非零多项式 \(g_1,\ldots,g_s\in R\)，可以经过有限次消项求出 \(q_1,\ldots,q_s,r\in R\)，使
\begin{equation}\label{b1:eq:F-division-representation}
f=q_1g_1+\cdots+q_sg_s+r,
\end{equation}
且 \(r\) 中的每个单项式都不能被任何 \(\operatorname{LM}(g_i)\) 整除。若 \(f\ne0\)，输出还满足：每个非零的 \(q_i g_i\) 以及非零的 \(r\)，其首单项式都不大于 \(\operatorname{LM}(f)\)。若 \(f=0\)，可取全部输出为零；若除式组为空，则取 \(r=f\)。
\end{theorem}
\begin{proof}
从 \(p=f\)、\(q_1=\cdots=q_s=r=0\) 开始，始终保持等式
\begin{equation}\label{b1:eq:F-division-invariant}
f=q_1g_1+\cdots+q_sg_s+r+p.
\end{equation}
只要 \(p\ne0\)，就按下列步骤处理。
\begin{enumerate}
\item 若有 \(i\) 满足 \(\operatorname{LM}(g_i)\mid\operatorname{LM}(p)\)，取符合条件的最小指标 \(i\)，令 \(t=\operatorname{LT}(p)/\operatorname{LT}(g_i)\)，把 \(q_i\) 改为 \(q_i+t\)，把 \(p\) 改为 \(p-tg_i\)。首项被消去；由\cref{b1:prop:F-leading-product}，其余留下的单项式均严格小于原来的 \(\operatorname{LM}(p)\)。
\item 若没有这样的 \(i\)，把 \(r\) 改为 \(r+\operatorname{LT}(p)\)，把 \(p\) 改为 \(p-\operatorname{LT}(p)\)。这一项无法再用任何除式的首项消去，所以移入余式。新的非零 \(p\) 的首单项式仍严格减小。
\end{enumerate}
两种操作都保持\eqref{b1:eq:F-division-invariant}。若过程永不结束，各个非零 \(p\) 的首单项式就形成无穷严格下降链，与单项式序的良序性矛盾。因此有限步后 \(p=0\)，得到\eqref{b1:eq:F-division-representation}。每个移入 \(r\) 的单项式都不被任何 \(\operatorname{LM}(g_i)\) 整除；相加至多消去已有的项，不会产生新的单项式，所以最终 \(r\) 仍满足这一条件。

最后检查首项界。每次加入某个 \(q_i\) 的项 \(t\) 都满足 \(\operatorname{LM}(tg_i)=\operatorname{LM}(p)\preceq\operatorname{LM}(f)\)，而移入 \(r\) 的项也满足同一上界。把这些项分别相加不会产生更大的单项式；由\cref{b1:prop:F-leading-product}，所有非零的 \(q_i g_i\) 和非零的 \(r\) 都满足所述上界。
\end{proof}

算法可以选择任意一个能够消项的除式，终止性和所述首项界都不变。规定“最小指标”只是把选择固定下来，使同一有序输入产生确定的输出。域的假设用在首项系数求逆处；良序性负责终止，而乘法相容性保证一次消项不会反而产生更大的单项式。这三件事在论证中各有不同作用。

\subsection{余式为何依赖除式顺序}
\label{b1:subsec:F01:03}

在 \(\mathbb Q[x,y]\) 中固定字典序 \(x\succ y\)，取
\[
g_1=xy-1,\qquad g_2=y^2-1,\qquad f=xy^2.
\]
\(f\) 的首单项式同时被 \(\operatorname{LM}(g_1)=xy\) 和 \(\operatorname{LM}(g_2)=y^2\) 整除。如果先用 \(g_1\)，就有
\[
xy^2-y(xy-1)=y,
\qquad f=yg_1+y.
\]
\(y\) 不能再被 \(xy\) 或 \(y^2\) 整除，所以余式为 \(y\)。如果先用 \(g_2\)，则
\[
xy^2-x(y^2-1)=x,
\qquad f=xg_2+x,
\]
这次余式为 \(x\)。两次计算均符合\cref{b1:thm:F-division}，但结果不同。差异不是计算失误，也不是系数近似所致；它来自允许消项的先后选择。

\ProseParagraphBreak
将两条表示相减，得到
\begin{equation}\label{b1:eq:F-missing-generator}
x-y=yg_1-xg_2\in(g_1,g_2).
\end{equation}
然而 \(x-y\) 的两项都不能被 \(xy\) 或 \(y^2\) 整除。因此用这两个除式去除 \(x-y\)，算法不消去任何一项，得到的非零余式就是 \(x-y\) 自身。对于任意给定的生成组，“余式非零”因而不能直接推出被除式不在所生成的理想中。

\ProseParagraphBreak
零余式的意义则没有这种不确定性。如果一次除法得到 \(r=0\)，\eqref{b1:eq:F-division-representation}就是 \(f\in(g_1,\ldots,g_s)\) 的具体证据，并记录了所需的系数 \(q_i\)。真正缺少的是反向判据：当 \(f\) 属于理想时，能否保证除法必得零余式？例子中的 \(x-y\) 提示了改进办法：原生成组没有直接显露理想中每一种可能的首项；把必要的关系补进去，才可能使消项过程准确反映整个理想。下一节把这一要求写成精确定义。

% !TeX root = ../../Book1.tex
% 纲要标题：### F.2 Gröbner 基与正规式
% 纲要位置：工作记录/计划纲要.md，第 975 行。

% 纲要要点（供目录核对）：
% #### F.2.1 单项式理想与有限性
% * 定义单项式理想，证明单项式属于它当且仅当被某个单项式生成元整除。
% * 用第一卷已证的 Hilbert 基定理推出有限生成性，说明它与 \(\mathbb N^n\) 按分量比较的 Dickson 有限性之间的对应。
% * 明确有限存在性与从给定多项式生成组实际求出标准生成组是两个步骤。
% #### F.2.2 首项理想与 Gröbner 基
% * 定义理想的首项理想以及 Gröbner 基。
% * 区分“给定生成元的首项所生成的理想”与“原理想中全部元素的首项所生成的理想”。
% * 从首项理想的有限单项式生成组提升出多项式，证明每个理想都存在有限 Gröbner 基，并证明这些多项式确实生成原理想。
% #### F.2.3 正规余式与理想判定
% * 证明对 Gröbner 基得到的不可约余式唯一；明确固定的是系数域和单项式序。
% * 证明理想成员判据：\(f\in I\) 当且仅当正规余式为零。
% * 用互相约化生成元检验理想包含和理想相等；保留零余式对应的线性组合等式。

\section{Gröbner 基与正规式}
\label{b1:sec:F02}

多元除法的每一步只检查单项式能否整除，因而改进除法首先要弄清一种特别简单的理想：由单项式生成的理想。它们的成员判定完全由指数决定。把一般理想中所有非零多项式的首单项式收集起来，再生成一个这样的理想，便能把“哪些首项必须能够消去”表达为一个有限条件。

\subsection{单项式理想与有限性}
\label{b1:subsec:F02:01}

\begin{definition}\label{b1:def:F-monomial-ideal}
设 \(k\) 为域，\(n\geq1\) 为整数，\(R=k[x_1,\ldots,x_n]\)。若理想 \(J\subseteq R\) 能由一族单项式生成，即
\[
J=(x^\alpha\mid\alpha\in A)
\]
对某个 \(A\subseteq\mathbb N^n\) 成立，则称 \(J\) 为\term[danxiangshi-lixiang]{单项式理想}[monomial ideal]。允许 \(A=\varnothing\)，此时 \(J=0\)。
\end{definition}

例如 \((x^2,xy,y^3)\subseteq k[x,y]\) 是单项式理想，而单个多项式 \(x-y\) 生成的理想不是：若它是单项式理想，下面的判据会使 \(x\) 和 \(y\) 都属于该理想，但把 \(y=x\) 代入 \(x=h(x,y)(x-y)\) 就会得到不成立的多项式等式 \(x=0\)。

\begin{proposition}\label{b1:prop:F-monomial-membership}
设 \(k\) 为域，\(n\geq1\) 为整数，\(R=k[x_1,\ldots,x_n]\)，\(J=(x^\alpha\mid\alpha\in A)\subseteq R\) 是单项式理想。单项式 \(x^\beta\) 属于 \(J\)，当且仅当它被某个 \(x^\alpha\)，其中 \(\alpha\in A\)，整除。更一般地，多项式 \(f\in R\) 属于 \(J\)，当且仅当 \(f\) 中的每个单项式都被某个这样的生成元整除。
\end{proposition}
\begin{proof}
若 \(f\in J\)，按生成理想的定义，存在有限个 \(\alpha^{(1)},\ldots,\alpha^{(m)}\in A\) 及多项式 \(h_1,\ldots,h_m\)，使
\[
f=\sum_{j=1}^m h_jx^{\alpha^{(j)}}.
\]
展开每个 \(h_j\)，右边出现的每个单项式都被对应的 \(x^{\alpha^{(j)}}\) 整除。合并同类项可能消去一些单项式，却不会引入从未出现过的单项式。因此 \(f\) 的每个单项式都被至少一个生成元整除。反过来，若 \(f\) 的每个单项式都如此，逐项选一个整除它的生成元，就把该项写成这个生成元的单项式倍数再乘一个系数。由于 \(f\) 只有有限项，把这些等式相加即得 \(f\in J\)。对 \(f=x^\beta\) 应用同一结论，便得到单项式的判据。
\end{proof}

这也说明单项式理想的成员判定不依赖单项式序。例如对 \(J=(x^2,xy,y^3)\)，\(x^4y+x y^2-y^5\) 的每一项都被列出的某个生成元整除，所以它在 \(J\) 中；\(x+y^2\) 的两项都不满足这一要求，故不在 \(J\) 中。对于单项式理想，不存在通过不同项相消而“意外产生”一个不被任何生成元整除的单项式的情况。

\begin{lemma}[单项式生成组的有限性]\label{b1:lem:F-monomial-finite}
设 \(k\) 为域，\(n\geq1\) 为整数，\(R=k[x_1,\ldots,x_n]\)。任一由 \(R\) 中单项式组成的集合 \(M\) 都有有限子集 \(M_0\subseteq M\)，使每个 \(u\in M\) 都被 \(M_0\) 中的某个单项式整除。因此单项式理想可以由有限多个单项式生成，而且可以从任何给定的单项式生成族中选取这些生成元。空集合取 \(M_0=\varnothing\)。
\end{lemma}
\begin{proof}
令 \(J=(M)\)。域只有零理想和单位理想，因而是 Noether 环。由\cref{b1:thm:hilbert-basis}，有限元多项式环 \(R=k[x_1,\ldots,x_n]\) 的每个理想有限生成，故有 \(J=(f_1,\ldots,f_t)\)。每个 \(f_i\) 都是 \(M\) 中有限多个单项式的多项式系数组合；把这有限多次出现的单项式收集起来，得到有限集 \(M_0\subseteq M\)，并有
\[
(M_0)\subseteq J=(f_1,\ldots,f_t)\subseteq(M_0).
\]
于是 \(J=(M_0)\)。任意 \(u\in M\) 属于这个理想，\cref{b1:prop:F-monomial-membership}就说明 \(u\) 被 \(M_0\) 中某个单项式整除。
\end{proof}

若把 \(\mathbb N^n\) 按分量比较，\(\alpha\leq\beta\) 表示对每个 \(i\) 都有 \(\alpha_i\leq\beta_i\)，那么引理等价于：每个集合 \(A\subseteq\mathbb N^n\) 都有有限子集 \(A_0\)，使每个 \(\beta\in A\) 都大于或等于某个 \(\alpha\in A_0\)。这称为\term[Dickson-yinli]{Dickson 引理}[Dickson's lemma]。从 \(A_0\) 中删去被其他成员严格压低的元素，可得 \(A\) 的全部极小元；它们仍为有限集。这里比较的是分量偏序，而不是为除法选定的字典序等全序。

\ProseParagraphBreak
还可以换一个角度看这一有限性。每当向一个单项式生成组中加入不被此前任何生成元整除的新单项式，\cref{b1:prop:F-monomial-membership}就保证相应单项式理想严格增大。\cref{b1:thm:hilbert-basis}和\cref{b1:prop:noetherian-acc}合起来说明，\(R\) 中的理想严格升链不可能无穷延续。下一节的算法正会不断加入这种“此前没有覆盖到”的首单项式。

\subsection{首项理想与 Gröbner 基}
\label{b1:subsec:F02:02}

回到一般理想 \(I\subseteq R\)，并重新固定一个单项式序 \(\prec\)。对任意给定生成组 \(g_1,\ldots,g_s\)，已经知道的首单项式只来自这几个元素；理想的其他元素却可能通过最高项相消产生新的首单项式。为了保证每个理想成员都能开始消项，必须考虑后者的全体。

\begin{definition}\label{b1:def:F-initial-ideal}
设 \(k\) 为域，\(n\geq1\) 为整数，\(R=k[x_1,\ldots,x_n]\) 配备单项式序 \(\prec\)，\(I\subseteq R\) 为理想。令
\[
\operatorname{in}_{\prec}(I)
\coloneqq\Bigl(\operatorname{LM}(f)\mid f\in I,\ f\ne0\Bigr)
\]
称 \(\operatorname{in}_{\prec}(I)\) 为 \(I\) 关于 \(\prec\) 的\term[shouxiang-lixiang]{首项理想}[initial ideal]。它是由 \(I\) 中所有非零多项式的首单项式生成的单项式理想。若 \(I=0\)，规定 \(\operatorname{in}_{\prec}(I)=0\)。
\end{definition}

\symidx{\(\operatorname{in}_{\prec}(I)\)：理想关于所选单项式序的首项理想}

因为每个非零系数在 \(k\) 中可逆，用所有首项 \(\operatorname{LT}(f)\) 生成也得到同一理想。因此“首项理想”的名称与这里使用首单项式的写法相容。这个定义依赖整个理想及所选次序，不依赖写下理想时采用了哪一组生成元。

\begin{definition}\label{b1:def:F-groebner-basis}
设 \(k\) 为域，\(n\geq1\) 为整数，\(R=k[x_1,\ldots,x_n]\) 配备单项式序 \(\prec\)，\(I\subseteq R\) 为理想，\(G=\{g_1,\ldots,g_s\}\) 为 \(I\) 中有限个非零多项式组成的集合。若
\begin{equation}\label{b1:eq:F-groebner-definition}
\operatorname{in}_{\prec}(I)
=\Bigl(\operatorname{LM}(g_1),\ldots,\operatorname{LM}(g_s)\Bigr),
\end{equation}
则称 \(G\) 为 \(I\) 关于 \(\prec\) 的一组\term[Groebner-ji]{Gröbner 基}[Gröbner basis]。零理想采用空集为 Gröbner 基。
\end{definition}

由\cref{b1:prop:F-monomial-membership}，\eqref{b1:eq:F-groebner-definition}等价于：对每个非零 \(f\in I\)，总有 \(g_i\in G\) 使 \(\operatorname{LM}(g_i)\mid\operatorname{LM}(f)\)。右边生成的单项式理想本来就包含于左边，所以只须保证每个理想成员的首单项式都被覆盖。这个定义暂时没有另行要求 \(G\) 生成 \(I\)；下面会证明生成性由所列条件推出。名称中的“基”也不表示这些多项式在线性代数意义下线性无关，更不表示每个理想成员的多项式系数表示唯一。

\ProseParagraphBreak
例如\eqref{b1:eq:F-missing-generator}说明，在字典序 \(x\succ y\) 下，\(I=(xy-1,y^2-1)\) 含有 \(x-y\)，所以 \(x\in\operatorname{in}_{\prec}(I)\)。但是 \(x\notin(xy,y^2)\)，因为它不被 \(xy\) 或 \(y^2\) 整除。因此原来的两个生成元不是 Gröbner 基。另一方面，任一单项式理想的有限单项式生成组就是 Gröbner 基：\cref{b1:prop:F-monomial-membership}已说明，该理想中每个多项式的首单项式都被某个生成元整除。

\begin{theorem}[Gröbner 基的存在性]\label{b1:thm:F-groebner-existence}
设 \(k\) 为域、\(n\geq1\) 为整数，\(R=k[x_1,\ldots,x_n]\)。对 \(R\) 的任意理想 \(I\) 及任意固定的单项式序，\(I\) 都存在有限 Gröbner 基。任一这样的基 \(G\) 都满足 \(I=(G)\)。
\end{theorem}
\begin{proof}
若 \(I=0\)，取空基即可。设 \(I\ne0\)，把\cref{b1:lem:F-monomial-finite}用于
\[
M=\Bigl\{\operatorname{LM}(f)\mid f\in I,\ f\ne0\Bigr\}.
\]
可从中选出有限个单项式生成整个 \((M)\)。按 \(M\) 的定义，为每个被选出的单项式取一个给出它的 \(g_i\in I\)。所得有限集 \(G=\{g_1,\ldots,g_s\}\) 满足\eqref{b1:eq:F-groebner-definition}，故为 Gröbner 基。

还须证明任何满足定义的 \(G\) 都生成 \(I\)。显然 \((G)\subseteq I\)。对 \(f\in I\)，按任意顺序用 \(G\) 作\cref{b1:thm:F-division}的除法，得到 \(f=\sum_iq_i g_i+r\)。于是 \(r\in I\)。若 \(r\ne0\)，由 Gröbner 基的定义及\cref{b1:prop:F-monomial-membership}，\(\operatorname{LM}(r)\) 被某个 \(\operatorname{LM}(g_i)\) 整除；这与除法余式中没有这样的单项式矛盾。因此 \(r=0\)，从而 \(f\in(G)\)，即 \(I\subseteq(G)\)。
\end{proof}

单位理想也没有例外：\(\{1\}\) 是 \(R\) 的 Gröbner 基，用它除任何多项式都得零余式。反过来，单位理想的任何 Gröbner 基必含非零常数，因为其首项理想含有 \(1\)，而唯一能够整除 \(1\) 的单项式就是 \(1\)。存在性定理尚未给出从一组任意多项式生成元算出 Gröbner 基的步骤；它说明所需的补充可以有限，但没有告诉我们应当依次补入哪些多项式。

\subsection{正规余式与理想判定}
\label{b1:subsec:F02:03}

用一组 Gröbner 基作除法时，商 \(q_i\) 仍可能不唯一。不过，余式所代表的是同一个剩余类，而且它被限制在同一组允许出现的单项式之内。这两条约束合在一起，足以保证余式唯一。

\begin{theorem}[正规式与成员判据]\label{b1:thm:F-normal-form}
设 \(k\) 为域，\(n\geq1\) 为整数，\(R=k[x_1,\ldots,x_n]\) 配备一个单项式序，\(I\subseteq R\) 为理想，\(G=\{g_1,\ldots,g_s\}\) 为 \(I\) 的 Gröbner 基。对任意 \(f\in R\)，存在唯一的 \(r\in R\)，同时满足
\begin{enumerate}
\item \(f-r\in I\)；
\item \(r\) 的每个单项式都不被任何 \(\operatorname{LM}(g_i)\) 整除。
\end{enumerate}
用 \(G\) 作除法，无论采用怎样的除式顺序或合法消项选择，余式都等于这个 \(r\)。称它为 \(f\) 关于 \(G\) 的\term[zhenggui-yushi]{正规余式}[normal remainder]，也称\term[zhenggui-shi]{正规式}[normal form]，记为 \(\operatorname{NF}_G(f)\)。并且 \(f\in I\) 当且仅当 \(\operatorname{NF}_G(f)=0\)。
\end{theorem}
\begin{proof}
\cref{b1:thm:F-division}给出至少一个这样的余式，因为 \(G\subseteq I\)，表示式 \(f- r=\sum_iq_i g_i\) 保证第一条。设 \(r\) 和 \(r'\) 都满足两条条件，则 \(r-r'\in I\)。两式相减只可能消去原有单项式，故 \(r-r'\) 中也没有被任何 \(\operatorname{LM}(g_i)\) 整除的单项式。若 \(r-r'\ne0\)，Gröbner 基的定义却要求它的首单项式被某个 \(\operatorname{LM}(g_i)\) 整除，矛盾。因此 \(r=r'\)。所有合法除法的余式均满足这两条，所以都相同。

若 \(r=0\)，第一条直接给出 \(f\in I\)。若 \(f\in I\)，零多项式已经满足两条条件，余式的唯一性便给出 \(r=0\)。
\end{proof}

\symidx{\(\operatorname{NF}_G(f)\)：关于 Gröbner 基的正规余式}
这里用 \(G\) 作下标是为了说明用于计算的基；实际上，允许出现的单项式恰是不属于 \(\operatorname{in}_{\prec}(I)\) 的单项式。因此在理想 \(I\) 和单项式序固定时，换用另一组 Gröbner 基也得到同一个正规余式。对一般的生成组，本附录只称“某次除法的余式”，不把它误记为唯一确定的 \(\operatorname{NF}_G(f)\)。

\begin{proposition}\label{b1:prop:F-normal-form-linear}
设 \(k\) 为域，\(n\geq1\) 为整数，\(R=k[x_1,\ldots,x_n]\) 配备一个单项式序，\(I\subseteq R\) 为理想，\(G\) 为 \(I\) 的一组 Gröbner 基。对 \(f,h\in R\) 和 \(a,b\in k\)，有
\[
\operatorname{NF}_G(af+bh)
=a\operatorname{NF}_G(f)+b\operatorname{NF}_G(h).
\]
此外，\(f-h\in I\) 当且仅当 \(\operatorname{NF}_G(f)=\operatorname{NF}_G(h)\)。
\end{proposition}
\begin{proof}
右边只含不属于 \(\operatorname{in}_{\prec}(I)\) 的单项式，而且它与 \(af+bh\) 的差属于 \(I\)。由\cref{b1:thm:F-normal-form}的唯一性，右边正是所求正规余式。取 \(a=1\)、\(b=-1\)，再用同一定理的成员判据，便得第二个断言。
\end{proof}

正规式通常不保持多项式的乘法。例如在 \(k[y]\) 中对 \(G=\{y^2\}\)，有 \(\operatorname{NF}_G(y)=y\)，但 \(\operatorname{NF}_G(y^2)=0\ne y^2\)。商环中的乘法需要先把代表元相乘，再对结果取正规式；\secref{b1:sec:F05}将据此完成具体计算。

\begin{proposition}[理想包含与相等的判定]\label{b1:prop:F-ideal-comparison}
设 \(k\) 为域，\(n\geq1\) 为整数，\(R=k[x_1,\ldots,x_n]\) 配备一个单项式序，\(I=(f_1,\ldots,f_t)\)、\(J=(h_1,\ldots,h_u)\) 为 \(R\) 的理想，并设 \(H\) 为 \(J\) 的一组 Gröbner 基。则
\[
I\subseteq J
\quad\Longleftrightarrow\quad
\operatorname{NF}_H(f_i)=0\quad(1\leq i\leq t).
\]
若还已知 \(I\) 的一组 Gröbner 基 \(G\)，则 \(I=J\) 当且仅当上述条件及 \(\operatorname{NF}_G(h_j)=0\)，其中 \(1\leq j\leq u\)，同时成立。
\end{proposition}
\begin{proof}
由\cref{b1:thm:F-normal-form}，所列余式为零恰好表示各个 \(f_i\) 都在 \(J\) 中。因为 \(J\) 是理想，这又等价于它包含这些元素的全部多项式系数组合，即包含 \(I\)。交换 \(I\) 和 \(J\) 得到反向包含的判据；两个包含同时成立便是相等。
\end{proof}

这里用于检验的 \(f_i\) 或 \(h_j\) 可以只是普通生成元，不要求它们本身已经是 Gröbner 基；承担消项的那一组则须有 Gröbner 基的保证。每个零余式还应保留除法给出的多项式组合，从而不仅得到“属于”或“不属于”的结论，也能核验相应的包含关系。若某次计算得到非零正规余式，则成员判据已经证明不存在把原多项式写成该理想生成元组合的表示。

% !TeX root = ../../Book1.tex
% 纲要标题：### F.3 Buchberger 判据与算法
% 纲要位置：工作记录/计划纲要.md，第 995 行。

% 纲要要点（供目录核对）：
% #### F.3.1 S 多项式与首项抵消
% * 利用首单项式的最小公倍式定义 S 多项式，说明其中的系数归一化。
% * 从两个最高项发生抵消的现象说明为何检查 S 多项式。
% * 完整陈述并证明 Buchberger 判据，不把正确性留给第三卷。
% #### F.3.2 逐步补充生成元的算法
% * 对所有待处理多项式对计算 S 多项式；把非零余式加入生成组，并登记新产生的待处理对。
% * 证明新增余式不改变原理想，却使首项理想严格增大；用理想升链条件证明算法终止。
% * 接续 F.1 的例子，算出 \((xy-1,y^2-1)\) 的 Gröbner 基 \(\{x-y,y^2-1\}\)，核验前后生成的理想相同。
% #### F.3.3 约化 Gröbner 基
% * 定义首一、首项冗余的删除及尾项约化。
% * 给出从 Gröbner 基得到约化基的步骤。
% * 证明固定系数域和单项式序下约化 Gröbner 基的唯一性；以换序算例说明条件不能遗漏。

\section{Buchberger 判据与算法}
\label{b1:sec:F03}

Gröbner 基的定义要求控制理想中所有非零多项式的首项，不能直接拿来作有限次检验。困难出现在相加时：几个多项式的最高项可能恰好抵消，留下一个原有首项没有控制住的新首项。本节把这种抵消归结到有限个多项式对，并据此给出从任意有限生成组出发的计算方法。始终固定域 \(k\)、多项式环 \(R=k[x_1,\ldots,x_n]\) 及一个全局单项式序 \(\prec\)。

\subsection{S 多项式与首项抵消}
\label{b1:subsec:F03:01}

先看非零多项式 \(f,g\)。为了消去它们的最高项，需要先把两个首单项式乘到同一个单项式；最小公倍式恰好给出所需的最小公共倍数。此外，还必须除去各自的首项系数，才能让两个被减的项具有相同系数。

\begin{definition}\label{b1:def:F-s-polynomial}
设 \(k\) 为域，\(n\geq1\) 为整数，\(R=k[x_1,\ldots,x_n]\) 配备一个单项式序，\(f,g\in R\) 均非零。记
\[
L=\operatorname{lcm}\Bigl(\operatorname{LM}(f),\operatorname{LM}(g)\Bigr)
\]
称
\begin{equation}\label{b1:eq:F-s-polynomial}
S(f,g)=
\frac{L}{\operatorname{LM}(f)}\frac{f}{\operatorname{LC}(f)}
-\frac{L}{\operatorname{LM}(g)}\frac{g}{\operatorname{LC}(g)}
\end{equation}
为 \(f,g\) 的 \term[S-duoxiangshi]{\(S\) 多项式}[S-polynomial]。这里单项式的商是按指数相减得到的单项式，首项系数的倒数在域 \(k\) 中计算。
\end{definition}
\symidx{\(S(f,g)\)：两个多项式的 S 多项式}

等式右侧两项的首项都是 \(L\)，所以 \(S(f,g)=0\)，或者 \(\operatorname{LM}\Bigl(S(f,g)\Bigr)\prec L\)。例如在 \(\mathbb Q[x,y]\) 中取字典序 \(x\succ y\)，对
\[
g_1=xy-1,\qquad g_2=y^2-1,
\]
有 \(L=xy^2\)，故
\begin{equation}\label{b1:eq:F-first-s-polynomial}
S(g_1,g_2)=y(xy-1)-x(y^2-1)=x-y.
\end{equation}
\(x-y\) 属于 \((g_1,g_2)\)，但其首单项式 \(x\) 不能被 \(xy\) 或 \(y^2\) 整除。原生成组因而不是 Gröbner 基。\subsecref{b1:subsec:F01:03}中两个不同余式 \(x\) 与 \(y\) 的差，正是这里新发现的关系。

\ProseParagraphBreak
下面的判据说明，不会再有一种必须单独检查的“多个最高项同时抵消”：任意有限多个最高项的抵消，都可以拆成两两之差。证明中保留这些差的首项界，才能保证替换之后确实比原来的表示更简单。

\begin{theorem}[Buchberger 判据]\label{b1:thm:F-buchberger-criterion}
设 \(k\) 为域，\(n\geq1\) 为整数，\(R=k[x_1,\ldots,x_n]\) 配备单项式序 \(\prec\)，\(G=(g_1,\ldots,g_s)\) 是 \(R\) 中非零多项式组成的有限有序组，并令 \(I=(g_1,\ldots,g_s)\)。下列条件等价：
\begin{equivalences}
\item\label{b1:item:F-buchberger-gb} \(G\) 是 \(I\) 关于 \(\prec\) 的 Gröbner 基；
\item\label{b1:item:F-buchberger-remainder} 对每个 \(i<j\)，将 \(S(g_i,g_j)\) 按此有序组作多元除法，所得余式为零；
\item\label{b1:item:F-buchberger-bound} 对每个 \(i<j\)，记
\(L_{ij}=\operatorname{lcm}\Bigl(\operatorname{LM}(g_i),\operatorname{LM}(g_j)\Bigr)\)，存在表示
\begin{equation}\label{b1:eq:F-s-bounded-representation}
S(g_i,g_j)=\sum_{\ell=1}^s h_{ij\ell}g_\ell,
\qquad
\operatorname{LM}(h_{ij\ell}g_\ell)\prec L_{ij}
\quad\text{若 }h_{ij\ell}\ne0.
\end{equation}
\end{equivalences}
允许 \(s=0\)；此时 \(I=0\)，后两项均无待检验的多项式对。
\end{theorem}

\begin{proof}
若\itemref{b1:item:F-buchberger-gb}成立，\(S(g_i,g_j)\in I\)，由\cref{b1:thm:F-normal-form}的理想成员判据，其除法余式为零，得到\itemref{b1:item:F-buchberger-remainder}。

若\itemref{b1:item:F-buchberger-remainder}成立，\cref{b1:thm:F-division}给出
\(S(g_i,g_j)=\sum_\ell h_{ij\ell}g_\ell\)。当 \(S(g_i,g_j)\ne0\) 时，同一定理给出各个非零乘积的首项界
\[
\operatorname{LM}(h_{ij\ell}g_\ell)
\preceq\operatorname{LM}\Bigl(S(g_i,g_j)\Bigr)
\prec L_{ij}.
\]
若 \(S(g_i,g_j)=0\)，取全部 \(h_{ij\ell}=0\) 即可。这证明\itemref{b1:item:F-buchberger-bound}。

现在假定\itemref{b1:item:F-buchberger-bound}。要证明 \(G\) 是 Gröbner 基，只须对每个非零 \(f\in I\)，证明某个 \(\operatorname{LM}(g_i)\) 整除 \(\operatorname{LM}(f)\)。把 \(g_i\) 首一化，记 \(\widetilde g_i=g_i/\operatorname{LC}(g_i)\)。展开 \(f=\sum_i q_i g_i\) 中各个系数多项式的项，便得到某种有限表示
\begin{equation}\label{b1:eq:F-buchberger-minimal-representation}
f=\sum_{\nu=1}^t a_\nu m_\nu\widetilde g_{i_\nu},
\qquad a_\nu\in k^\times,\quad m_\nu\text{ 为单项式}.
\end{equation}
对每种这样的表示，取其中出现的最大首单项式
\[
M=\max_\nu\Bigl\{m_\nu\operatorname{LM}(g_{i_\nu})\Bigr\}.
\]
所有可能的 \(M\) 构成非空单项式集合；由 \(\prec\) 的良序性，可以在全部表示中选取 \(M\) 最小的一种。相加不能产生大于各被加式首单项式的项，所以 \(\operatorname{LM}(f)\preceq M\)。如果取等号，任取达到 \(M\) 的一项，便有 \(\operatorname{LM}(g_{i_\nu})\mid\operatorname{LM}(f)\)，结论已经成立。

若 \(\operatorname{LM}(f)\prec M\)，令
\(J=\Bigl\{\nu\mid m_\nu\operatorname{LM}(g_{i_\nu})=M\Bigr\}\)。由于 \(\widetilde g_i\) 首一，\(M\) 的系数抵消正是 \(\sum_{\nu\in J}a_\nu=0\)。固定 \(\nu_0\in J\)，得到
\begin{equation}\label{b1:eq:F-leading-cancellation}
\sum_{\nu\in J}a_\nu m_\nu\widetilde g_{i_\nu}
=\sum_{\substack{\nu\in J\\\nu\ne\nu_0}}
a_\nu\Bigl(m_\nu\widetilde g_{i_\nu}
-m_{\nu_0}\widetilde g_{i_{\nu_0}}\Bigr).
\end{equation}
若 \(i_\nu=i_{\nu_0}\)，则两个首单项式相同迫使 \(m_\nu=m_{\nu_0}\)，对应的差为零。否则令
\[
L=\operatorname{lcm}\Bigl(
\operatorname{LM}(g_{i_\nu}),
\operatorname{LM}(g_{i_{\nu_0}})\Bigr).
\]
两者都整除 \(M\)，故 \(L\mid M\)，且按\eqref{b1:eq:F-s-polynomial}，
\[
m_\nu\widetilde g_{i_\nu}
-m_{\nu_0}\widetilde g_{i_{\nu_0}}
=\frac{M}{L}S(g_{i_\nu},g_{i_{\nu_0}}).
\]
若指标次序相反，只须使用 \(S(g_j,g_i)=-S(g_i,g_j)\)。由\itemref{b1:item:F-buchberger-bound}，右侧可表示成若干 \(g_\ell\) 的多项式倍数，且每个非零乘积的首单项式严格小于 \((M/L)L=M\)；这里用了单项式序与乘法相容。把各个系数多项式再展开成项，不会使该首项界增大。

因此，\eqref{b1:eq:F-leading-cancellation}左侧的全部最高项，都能用首单项式严格小于 \(M\) 的项替代。原表示中不在 \(J\) 内的项也已具有这个性质。替换后便得到 \(f\) 的新表示，其最大首单项式严格小于 \(M\)，与所选的最小性矛盾。故只能有 \(\operatorname{LM}(f)=M\)，证明了\itemref{b1:item:F-buchberger-gb}。
\end{proof}

这个结果称为\term[Buchberger-panju]{Buchberger 判据}[Buchberger criterion]。实际检验通常使用余式为零的形式；带首项界的表示则解释了判据为什么成立。单说 \(S(g_i,g_j)\in I\) 没有检验作用，因为这对任何生成组都成立，必须控制它如何由给定生成元表示。

\subsection{逐步补充生成元的算法}
\label{b1:subsec:F03:02}

若某个 \(S\) 多项式留下非零余式，这个余式仍在原理想中，而它的首单项式尚未被已有生成元控制。把它加入生成组，恰好补充一条有用的新关系。新关系又会产生新的多项式对，所以不能只检查最初的那些对；应保存一个尚未处理的对的集合，直到其中没有元素。

\ProseParagraphBreak
\cref{b1:alg:F-buchberger}给出\term[Buchberger-suanfa]{Buchberger 算法}[Buchberger algorithm]的基本形式。系数域须支持精确的四则运算与零判定；有理数域和素域满足这一要求。算法不在中途删除生成元，因而各个下标始终保留原有含义。

\begin{algorithm}[htbp]
\caption{从有限生成组计算 Gröbner 基}
\label{b1:alg:F-buchberger}
\begin{algorithmsteps}{有限组 \(f_1,\ldots,f_m\in k[x_1,\ldots,x_n]\) 及全局单项式序 \(\prec\)。}{理想 \((f_1,\ldots,f_m)\) 的一个 Gröbner 基 \(G\)。}
\State 删除输入中的零多项式，将其余多项式首一化，依次记为 \(G=(g_1,\ldots,g_s)\)。
\State 令 \(P=\Bigl\{(i,j)\mid 1\le i<j\le s\Bigr\}\)。
\While{\(P\ne\varnothing\)}
\State 从 \(P\) 中取出一对 \((i,j)\)，并从 \(P\) 删除它。
\State 将 \(S(g_i,g_j)\) 按当前有序组 \(G\) 作多元除法，记余式为 \(r\)。
\If{\(r\ne0\)}
\State 令 \(g_{s+1}=r/\operatorname{LC}(r)\)，将其添到 \(G\) 末尾。
\State 把所有 \((1,s+1),\ldots,(s,s+1)\) 加入 \(P\)，再令 \(s\) 增加一。
\EndIf
\EndWhile
\State \Return \(G\)。
\end{algorithmsteps}
\end{algorithm}

\begin{theorem}\label{b1:thm:F-buchberger-algorithm}
设 \(k\) 为支持精确四则运算与零判定的域，\(n\geq1\) 为整数，\(R=k[x_1,\ldots,x_n]\) 配备全局单项式序 \(\prec\)，并假设任意两个单项式在 \(\prec\) 下的比较可在有限步内完成。对有限组 \(f_1,\ldots,f_m\in R\)，以这组多项式和 \(\prec\) 为输入的\cref{b1:alg:F-buchberger}在有限步后终止，输出的 \(G\) 是理想 \((f_1,\ldots,f_m)\) 关于 \(\prec\) 的 Gröbner 基。
\end{theorem}

\begin{proof}
记原理想为 \(I\)。每次除法给出等式
\begin{equation}\label{b1:eq:F-buchberger-step}
S(g_i,g_j)=\sum_{\ell=1}^s q_\ell g_\ell+r.
\end{equation}
左侧属于当前生成组所生成的理想，因而 \(r\) 也属于它。添加 \(r/\operatorname{LC}(r)\) 不改变所生成的理想，所以算法始终生成 \(I\)。

若 \(r\ne0\)，\cref{b1:thm:F-division}说明 \(r\) 的任何单项式都不能被当前的 \(\operatorname{LM}(g_\ell)\) 整除。由\cref{b1:prop:F-monomial-membership}，\(\operatorname{LM}(r)\) 不在它们生成的单项式理想内。故每添加一次非零余式，就有严格包含
\[
\Bigl(\operatorname{LM}(g_1),\ldots,\operatorname{LM}(g_s)\Bigr)
\subsetneq
\Bigl(\operatorname{LM}(g_1),\ldots,\operatorname{LM}(g_s),
\operatorname{LM}(r)\Bigr).
\]
由\cref{b1:thm:hilbert-basis}，\(R\) 是 Noether 环；再由\cref{b1:prop:noetherian-acc}，其中不可能存在无限严格上升的理想链。因此非零余式只会添加有限次。每次只引入有限个新的对，而每轮处理一个对且不将它重复放回，故全部对也在有限轮内处理完毕。每一轮的多元除法本身由\cref{b1:thm:F-division}保证终止，所以整个算法终止。

还须说明，早先处理过的对不会因生成组扩大而失去作用。处理 \((i,j)\) 时，若 \(S(g_i,g_j)=0\)，它已有全零的表示。否则，多元除法的首项界说明\eqref{b1:eq:F-buchberger-step}中每个非零 \(q_\ell g_\ell\)，以及可能非零的 \(r\)，其首单项式均不大于 \(\operatorname{LM}\Bigl(S(g_i,g_j)\Bigr)\)，因而严格小于 \(L_{ij}\)。如果 \(r=0\)，这直接给出\eqref{b1:eq:F-s-bounded-representation}；如果 \(r\ne0\)，把最后一项写成 \(\operatorname{LC}(r)g_{s+1}\)，也给出同样的表示。算法保留全部已有生成元，所以这份表示在以后的每一步仍然成立。

结束时每对输出生成元都已处理，因而全部满足\cref{b1:thm:F-buchberger-criterion}的\itemref{b1:item:F-buchberger-bound}。该判据说明输出确实是 \(I\) 的 Gröbner 基。若输入全为零，算法直接返回空组，也符合零理想的约定。
\end{proof}

终止性只排除了无限计算，并没有给出实用的运行时间上界。变量个数、次数、系数域和单项式序都会影响中间多项式的大小。初学时宜先把每一步的理想等式保存下来；它既便于查错，也能在最终判定 \(f\in I\) 时回代出 \(f\) 关于原生成元的表示。

\ProseParagraphBreak
继续计算\eqref{b1:eq:F-first-s-polynomial}中的例子。\(x-y\) 对 \((g_1,g_2)\) 已经不可约，所以添入
\[
g_3=x-y=yg_1-xg_2.
\]
新产生的两个多项式对满足
\begin{align*}
S(g_1,g_3)&=g_1-yg_3=y^2-1=g_2,\\
S(g_2,g_3)&=xg_2-y^2g_3=y^3-x=yg_2-g_3.
\end{align*}
它们的除法余式均为零，因此 \((g_1,g_2,g_3)\) 已是 Gröbner 基。其中 \(g_1\) 是冗余的，因为
\begin{equation}\label{b1:eq:F-example-ideal-certificates}
g_3=yg_1-xg_2,\qquad g_1=yg_3+g_2.
\end{equation}
这两个等式分别证明加入 \(g_3\) 和删去 \(g_1\) 不改变理想，故
\[
(xy-1,y^2-1)=(x-y,y^2-1).
\]
两个保留生成元的首单项式为 \(x,y^2\)，仍然生成原来的首项理想。因此
\begin{equation}\label{b1:eq:F-main-reduced-basis}
G=\{x-y,y^2-1\},\qquad
\operatorname{in}_\prec(I)=(x,y^2)
\end{equation}
给出一组 Gröbner 基，而且它已经具有下一小节所要求的约化形式。

\subsection{约化 Gröbner 基}
\label{b1:subsec:F03:03}

算法处理多项式对的次序可以改变中间结果；即使已经得到 Gröbner 基，仍可添入理想中的其他多项式而不破坏这一性质。若希望同一个理想得到可以直接比较的输出，还需要删去多余的首项，并把每个生成元中可以由其他首项约去的项消掉。

\begin{definition}\label{b1:def:F-reduced-groebner}
设 \(k\) 为域，\(n\geq1\) 为整数，\(R=k[x_1,\ldots,x_n]\) 配备单项式序 \(\prec\)，\(I\subseteq R\) 为理想，\(G\subseteq I\setminus\{0\}\) 为有限 Gröbner 基。若每个 \(g\in G\) 都首一，且对任意不同的 \(g,h\in G\)，\(g\) 中的任何单项式都不能被 \(\operatorname{LM}(h)\) 整除，则称 \(G\) 为 \(I\) 关于 \(\prec\) 的\term[yuehua-Groebner-ji]{约化 Gröbner 基}[reduced Gröbner basis]。零理想的约化 Gröbner 基约定为空集。
\end{definition}

这个条件同时约束首项和其余项。它要求不同首单项式之间不能互相整除；也要求每个多项式的尾项对其他生成元已经不可约。仅将首项系数化成一，并不足以得到约化基。例如在字典序 \(x\succ y\) 下，\(\{x-y,y^2-1,xy-1\}\) 首一且为 Gröbner 基，但 \(xy-1\) 的首单项式仍被 \(x-y\) 的首单项式整除。

\begin{theorem}\label{b1:thm:F-reduced-groebner}
设 \(k\) 为域，\(n\geq1\) 为整数，\(R=k[x_1,\ldots,x_n]\) 配备全局单项式序 \(\prec\)。则 \(R\) 的每个理想都有唯一的关于 \(\prec\) 的约化 Gröbner 基。给定该理想的一个有限 Gröbner 基，可以通过删除首项冗余、首一化及多元除法得到这个约化基。
\end{theorem}

\begin{proof}
先证存在性。零理想使用空集，以下设 \(I\ne0\)。由\cref{b1:thm:F-groebner-existence}取一个有限 Gröbner 基。只要某个生成元的首单项式被另一个的首单项式整除，就删去前者；首单项式相同时也只保留其中一个。每次删除后，保留首单项式生成的理想不变，仍为 \(\operatorname{in}_\prec(I)\)，所以保留下来的多项式仍组成 \(I\) 的 Gröbner 基。由于生成元有限，这一步终止。再将每个保留生成元首一化，记为 \(g_1,\ldots,g_t\)，并记 \(m_i=\operatorname{LM}(g_i)\)。这些 \(m_i\) 两两互不整除。

对每个 \(i\)，把尾项多项式 \(g_i-m_i\) 除以 \(g_1,\ldots,g_{i-1},g_{i+1},\ldots,g_t\)，记余式为 \(r_i\)，并令
\[
h_i=m_i+r_i.
\]
这一步可对同一组原来的 \(g_j\) 分别进行，不要求边计算边替换除式。如果尾项非零，\cref{b1:thm:F-division}保证 \(r_i\) 中所有单项式均严格小于 \(m_i\)，且不能被任何 \(m_j\)、\(j\ne i\)，整除。尾项为零时取 \(r_i=0\)。由于一个被 \(m_i\) 整除的单项式不可能严格小于 \(m_i\)，\(r_i\) 也没有被 \(m_i\) 整除的项。

除法等式说明 \(g_i-h_i\) 是其他 \(g_j\) 的组合，故 \(h_i\in I\)。同时 \(h_i\) 首一且 \(\operatorname{LM}(h_i)=m_i\)。因此 \(H=\{h_1,\ldots,h_t\}\) 的首单项式仍生成 \(\operatorname{in}_\prec(I)\)，从而是 \(I\) 的 Gröbner 基；前一段对余式的说明又表明它约化。这就证明了存在性和所述计算步骤。

再证唯一性。设 \(G,H\) 都是 \(I\) 的约化 Gröbner 基。它们的首单项式分别生成同一个单项式理想 \(\operatorname{in}_\prec(I)\)，且在每一组内部两两互不整除。对任意 \(g\in G\)，由\cref{b1:prop:F-monomial-membership}，存在 \(h\in H\) 使 \(\operatorname{LM}(h)\mid\operatorname{LM}(g)\)；反过来又有 \(g'\in G\) 使 \(\operatorname{LM}(g')\mid\operatorname{LM}(h)\)。于是
\[
\operatorname{LM}(g')\mid\operatorname{LM}(h)\mid\operatorname{LM}(g).
\]
\(G\) 内部互不整除，迫使 \(g'=g\)，所以三个首单项式相等。交换 \(G,H\) 的角色，得到它们具有完全相同的一组首单项式。每个首单项式在各组内恰好出现一次，因此可以按首单项式将两组一一配对。

取配对的 \(g,h\)，其首单项式同为 \(m\)。两者均首一，所以 \(g-h\) 的最高项抵消；它剩下的每个可能非零单项式都来自 \(g\) 或 \(h\) 的尾项。约化性排除了这些尾项被其他首单项式整除，而它们严格小于 \(m\)，也不能被 \(m\) 整除。因此 \(g-h\) 中没有单项式属于 \(\operatorname{in}_\prec(I)\)。但 \(g-h\in I\)，若它非零，其首单项式按定义必须属于 \(\operatorname{in}_\prec(I)\)，矛盾。故 \(g=h\)，所有配对元素均相同，从而 \(G=H\)。
\end{proof}

单位理想的约化 Gröbner 基是 \(\{1\}\)。因为其首项理想含有单项式 \(1\)，约化基中必须有首单项式为 \(1\) 的首一多项式，而这样的多项式只能是 \(1\)；它的首单项式整除所有其他单项式，因而约化性不允许再有别的生成元。这样零理想和单位理想也都被唯一性结论包括在内。

\ProseParagraphBreak
固定单项式序的要求不能省略。对同一个理想
\(I=(xy-1,y^2-1)\subseteq\mathbb Q[x,y]\)，字典序 \(x\succ y\) 给出约化基 \(\{x-y,y^2-1\}\)。若改用字典序 \(y\succ x\)，约化基则为
\[
\{y-x,x^2-1\}.
\]
两组确实生成同一个理想，因为
\begin{align*}
y-x&=-(x-y),&
x^2-1&=(x+y)(x-y)+(y^2-1),\\
xy-1&=x(y-x)+(x^2-1),&
y^2-1&=(x+y)(y-x)+(x^2-1).
\end{align*}
新次序下两个首单项式为 \(y,x^2\)，其 \(S\) 多项式满足
\[
x^2(y-x)-y(x^2-1)
=y-x^3=(y-x)-x(x^2-1).
\]
按两个生成元依次除去后余式为零，所以由\cref{b1:thm:F-buchberger-criterion}得到 Gröbner 基；其尾项显然不可约，故也是约化基。改变次序并未改变原理想，但改变了哪些项应当保留为余式。这种选择在消去指定变量时尤其有用。

% !TeX root = ../../Book1.tex
% 纲要标题：### F.4 消去与多项式方程组
% 纲要位置：工作记录/计划纲要.md，第 1015 行。

% 纲要要点（供目录核对）：
% #### F.4.1 字典序与消去定理
% * 定义 \(I\cap k[x_{r+1},\ldots,x_n]\) 这样的消去理想。
% * 证明字典序下，从 Gröbner 基中取出不含被消去变量的多项式，得到消去理想的一组 Gröbner 基。
% * 强调结论依赖所选单项式序，不把任意 Gröbner 基的子集直接当作消去基。
% #### F.4.2 消去后的方程与原方程的解
% * 用二元方程组演示先求消去方程，再代回恢复其余变量；每一步说明所在系数域和所求解的范围。
% * 完成一个简单的参数消去算例。
% * 回引第一卷第12.4节的 \(xy-1=0\) 例子，说明满足消去后关系的点未必能提升为原方程的解。
% * 几何投影、Zariski 闭包及零点定理的系统解释归第三卷。

\section{消去与多项式方程组}
\label{b1:sec:F04}

把一组方程改写成容易逐次求解的形式，是计算 Gröbner 基的一项直接用途。理想不变保证全部原关系仍被保存；要让其中一部分关系不再含某些变量，还须有适当的单项式序。本节先证明代数上的消去结论，再区分消去方程的解与原方程的解。

\subsection{字典序与消去定理}
\label{b1:subsec:F04:01}

设 \(k\) 为域，\(n\geq1\) 为整数，\(R=k[x_1,\ldots,x_n]\)，\(I\subseteq R\) 为理想。对整数 \(0\leq r\leq n\)，令
\[
R_r=k[x_{r+1},\ldots,x_n],\qquad I_r=I\cap R_r,
\]
其中 \(R_n=k\)。理想 \(I_r\) 称为消去前 \(r\) 个变量的\term[xiaoqu-lixiang]{消去理想}[elimination ideal]。它收集全部只用剩余变量写出的原理想关系；并不是随意删除生成元中含有其他变量的项。
\symidx{\(I_r=I\cap k[x_{r+1},\ldots,x_n]\)：消去理想}

\begin{theorem}[消去定理]\label{b1:thm:F-elimination}
设 \(k\) 为域，\(n\geq1\) 为整数，\(R=k[x_1,\ldots,x_n]\)，\(I\subseteq R\) 为理想，\(G\) 是 \(I\) 关于字典序 \(x_1\succ\cdots\succ x_n\) 的 Gröbner 基。对整数 \(0\leq r\leq n\)，当 \(r<n\) 时令 \(R_r=k[x_{r+1},\ldots,x_n]\)，并令 \(R_n=k\)；再置
\[
I_r=I\cap R_r,\qquad G_r=G\cap R_r.
\]
则 \(G_r\) 是 \(I_r\) 关于限制字典序的 Gröbner 基。
\end{theorem}
\begin{proof}
字典序有如下性质：若一个多项式的首单项式不含 \(x_1,\ldots,x_r\)，则它的每个单项式都不含这些变量。否则取某个含这些变量的项，并看其中最早出现的变量；与只含剩余变量的首单项式比较，该项反而较大，矛盾。

现在取 \(0\ne f\in I_r\)。由\cref{b1:def:F-groebner-basis}及\cref{b1:prop:F-monomial-membership}，存在 \(g\in G\)，使 \(\operatorname{LM}(g)\mid\operatorname{LM}(f)\)。所以 \(\operatorname{LM}(g)\) 不含前 \(r\) 个变量，上述性质使整个 \(g\in R_r\)，即 \(g\in G_r\)。因此每个 \(I_r\) 中非零多项式的首单项式，都被 \(G_r\) 中某个首单项式整除。

还须说明 \(G_r\) 生成 \(I_r\)。在 \(R_r\) 中用 \(G_r\) 对 \(f\in I_r\) 作除法，余式 \(h\) 仍在 \(I_r\)。若 \(h\ne0\)，刚证的整除性质与余式中没有可约单项式矛盾。故余式为零，除法等式说明 \(f\in(G_r)\)。于是生成性与首项条件同时成立。若 \(I_r=0\)，空集就是相应的 Gröbner 基，论证也包括这一情形。
\end{proof}

这个结果称为\term[xiaoqu-dingli]{消去定理}[elimination theorem]。对前面反复计算的 \(I=(xy-1,y^2-1)\subseteq\mathbb Q[x,y]\)，式\eqref{b1:eq:F-main-reduced-basis}给出字典序 \(x\succ y\) 下的约化基 \(\{x-y,y^2-1\}\)。所以
\begin{equation}\label{b1:eq:F-basic-elimination}
I\cap\mathbb Q[y]=(y^2-1).
\end{equation}
要消去 \(y\) 而保留 \(x\)，应选取 \(y\succ x\) 的字典序，得到 \(\{y-x,x^2-1\}\)。变量次序说明了算法希望优先消去哪一部分，而不是方程原先书写的先后次序。

\subsection{消去后的方程与原方程的解}
\label{b1:subsec:F04:02}

在任意包含 \(k\) 的域 \(K\) 中，两组生成同一理想的多项式具有相同的公共零点。事实上，若新生成元是旧生成元的多项式线性组合，旧生成元同时为零时，新生成元也为零；用反向表示得到另一包含。这里不需要零点定理，只用已经记录的理想等式。

\begin{proposition}[三角生成组]\label{b1:prop:F-triangular-basis}
设 \(k\) 为域，\(R=k[x,y]\)，\(h,p\in k[y]\)，其中 \(p\) 首一且次数 \(d\ge1\)。对字典序 \(x\succ y\)，\(\Bigl\{x-h(y),p(y)\Bigr\}\) 是理想 \(I=\Bigl(x-h(y),p(y)\Bigr)\subseteq R\) 的 Gröbner 基。对任意 \(f\in R\)，其正规余式等于先算 \(f\Bigl(h(y),y\Bigr)\)，再除以 \(p(y)\) 所得的一元余式。
\end{proposition}
\begin{proof}
两个首单项式为 \(x,y^d\)。对 \(0\ne f\in I\)，若其首单项式含 \(x\)，就被第一个整除。否则由字典序的性质，\(f\in k[y]\)。在它的理想表示中代入 \(x=h(y)\)，得到 \(f(y)\) 是 \(p(y)\) 的倍数，故其首单项式被 \(y^d\) 整除。这个整除性质及单项式理想成员判据说明给定生成组是 Gröbner 基。

对每个 \(m\ge1\)，恒等式
\[
x^m-h(y)^m=\Bigl(x-h(y)\Bigr)\sum_{j=0}^{m-1}x^{m-1-j}h(y)^j
\]
说明 \(f(x,y)\) 与代入后的多项式模 \(x-h(y)\) 同余。再作一元除法，得到仅含 \(1,y,\ldots,y^{d-1}\) 的余式，它与 \(f\) 模 \(I\) 同余且不可约，由\cref{b1:thm:F-normal-form}必为正规余式。
\end{proof}

因此，对上面的有理数算例，先解 \(y^2-1=0\)，再从 \(x-y=0\) 得 \(x=y\)，得到全部有理解及复数解 \((1,1),(-1,-1)\)。这次代回没有丢失条件，因为手中仍保存着与原理想相同的整组 Gröbner 基，而不只有消去方程。

\begin{example}[参数的消去]\label{b1:exm:F-parameter-elimination}
考虑参数式 \(x=t,y=t^2\)，在 \(k[t,x,y]\) 中取
\[
J=(x-t,y-t^2),\qquad t\succ x\succ y.
\]
令 \(g_1=t-x\)、\(g_2=x^2-y\)。有
\[
g_1=-(x-t),\qquad g_2=(t+x)(x-t)-(y-t^2),
\]
以及 \(x-t=-g_1\)、\(y-t^2=-g_2-(t+x)g_1\)，所以 \((g_1,g_2)=J\)。其首单项式为 \(t,x^2\)，唯一需要检查的 S 多项式为
\[
x^2g_1-tg_2=ty-x^3=yg_1-xg_2.
\]
按 \(g_1,g_2\) 约化，先减去 \(yg_1\)，再消去 \(-xg_2\)，余式为零。\cref{b1:thm:F-buchberger-criterion}说明它们构成 Gröbner 基。因此
\[
J\cap k[x,y]=(x^2-y).
\]
这个例子中，反过来满足 \(y=x^2\) 的每个域值点都能提升：取 \(t=x\) 即可。该提升来自具体公式，并非一般消去定理的一部分。
\end{example}

\begin{example}[不能提升的点]\label{b1:exm:F-projection-nonlifting}
对 \(I=(xy-1)\subseteq k[x,y]\)，有 \(I\cap k[x]=0\)，但原方程没有第一坐标为零的解。
\end{example}
\begin{proof}
若 \(0\ne h\in k[x,y]\)，把多项式看成 \(k[x][y]\) 中的元素。因 \(k[x]\) 是整环，\((xy-1)h\) 关于 \(y\) 的次数为 \(1+\deg_yh\ge1\)，不可能属于 \(k[x]\)。故交理想为零。零理想对剩余坐标没有方程限制，可是 \(x=0\) 时 \(xy-1=-1\ne0\)，所以该点没有提升。
\end{proof}

这将\subsecref{b1:subsec:1204:01}中的现象放进了消去理想的框架。另一个容易混淆的问题是解所在的域。例如 \((x-y,y^2+1)\subseteq\mathbb Q[x,y]\) 的计算可以完全在有理数中进行，但它没有有理解，在复数域中才有 \((i,i),(-i,-i)\)。系数运算所在的域与允许取解的域应分别写明。

\ProseParagraphBreak
消去定理本身不要求基域代数闭。若进一步研究投影像与其 Zariski 闭包的关系，则需要第三卷第8.3节的几何解释；本附录只使用这里已经证明的理想等式和逐点代回。

% !TeX root = ../../Book1.tex
% 纲要标题：### F.5 标准单项式与商环中的计算
% 纲要位置：工作记录/计划纲要.md，第 1030 行。

% 纲要要点（供目录核对）：
% #### F.5.1 标准单项式构成商环的基
% * 定义不属于首项理想的标准单项式。
% * 用正规余式证明：这些单项式的剩余类构成 \(k[x_1,\ldots,x_n]/I\) 的一组 \(k\) 基，分别说明生成性与独立性。
% * 证明商环有限维当且仅当每个变量的某个正次幂属于首项理想，并用标准单项式计数维数。
% #### F.5.2 商环的运算、单位与零因子
% * 用“先相加或相乘，再取正规余式”的办法完成商环运算。
% * 在 \(\mathbb Q[x,y]/(x-y,y^2-1)\) 中使用基 \(1,y\) 写出乘法，并计算零因子、幂等元和可逆元。
% * 证明并运用单位判据：\(f+I\) 可逆当且仅当 \(1\in I+(f)\)；从生成关系中提取逆元。
% * 与 \(\mathbb Q[x,y]/(x,y^2)\) 比较，说明商环维数相同并不意味着不同解的个数相同。
% 节末用一个短段比较交换单项式和非交换字：前者按指数比较整除，后者需要保留字母顺序并检查连续子字的重叠。阅读第二卷第20.6节时，再系统学习 Diamond 引理及非交换 Gröbner 基。
% ---

\section{标准单项式与商环中的计算}
\label{b1:sec:F05}

正规余式不只回答一个多项式是否属于理想，还为商环的每个元素选出一个确定的坐标表达。这样，商环中的乘法可以回到多项式乘法，再用同一个 Gröbner 基整理。本节只使用向量空间的通常知识，不要求预先学习第二卷的一般模论。

\subsection{标准单项式构成商环的基}
\label{b1:subsec:F05:01}

设 \(k\) 为域，\(n\geq1\) 为整数，\(R=k[x_1,\ldots,x_n]\) 配备全局单项式序 \(\prec\)，\(I\subseteq R\) 为理想，\(G\) 为 \(I\) 的 Gröbner 基。若单项式不属于 \(\operatorname{in}_{\prec}(I)\)，则称它为 \(I\) 关于 \(\prec\) 的\term[biaozhun-danxiangshi]{标准单项式}[standard monomial]，全体记为 \(\mathcal B\)。由\cref{b1:prop:F-monomial-membership}，这恰好是不被任何 \(\operatorname{LM}(g)\)、\(g\in G\)，整除的单项式。
\symidx{\(\mathcal B\)：给定理想及单项式序下的标准单项式集合}

\begin{theorem}[标准单项式基]\label{b1:thm:F-standard-monomial-basis}
设 \(k\) 为域，\(n\geq1\) 为整数，\(R=k[x_1,\ldots,x_n]\) 配备单项式序 \(\prec\)，\(I\subseteq R\) 为理想，\(G\) 为 \(I\) 的 Gröbner 基。令 \(\mathcal B\) 为不属于 \(\operatorname{in}_{\prec}(I)\) 的全部单项式所成集合。则 \(\mathcal B\) 中单项式的剩余类构成 \(R/I\) 的一组 \(k\) 向量空间基；对任意 \(f\in R\)，商类 \(f+I\) 在这组基下的坐标由正规余式 \(\operatorname{NF}_G(f)\) 给出。
\end{theorem}
\begin{proof}
\cref{b1:thm:F-division}将 \(f\) 写成 \(\sum_iq_ig_i+r\)，其中 \(r\) 只含标准单项式。所以 \(f+I=r+I\)，证明这些剩余类生成商空间。

若标准单项式之间有非平凡的有限线性关系，其所代表的非零多项式 \(h\) 属于 \(I\)。但是 \(h\) 的首单项式仍是标准单项式，不能属于 \(\operatorname{in}_{\prec}(I)\)，与 \(h\in I\) 矛盾。因此它们线性无关。上述余式就是\cref{b1:thm:F-normal-form}的唯一正规余式，给出唯一的基坐标。
\end{proof}

\begin{corollary}\label{b1:cor:F-finite-quotient}
设 \(k\) 为域，\(n\geq1\) 为整数，\(R=k[x_1,\ldots,x_n]\) 配备单项式序 \(\prec\)，\(I\subseteq R\) 为理想，并令 \(\mathcal B\) 为相应的标准单项式集合。则 \(R/I\) 为有限维 \(k\) 空间，当且仅当对每个变量 \(x_i\)，存在正整数 \(d_i\) 使 \(x_i^{d_i}\in\operatorname{in}_{\prec}(I)\)。成立时
\[
\dim_k(R/I)=|\mathcal B|\le d_1\cdots d_n.
\]
\end{corollary}
\begin{proof}
若每个这样的幂都属于首项理想，则标准单项式的第 \(i\) 个指数小于 \(d_i\)，所以只可能有至多 \(d_1\cdots d_n\) 个。基定理便给出有限维性及计数公式。反过来，若对某个 \(i\) 的全部正次幂均不属于首项理想，\(1,x_i,x_i^2,\ldots\) 全为标准单项式，基定理使商空间无限维，矛盾。若 \(I=R\)，标准集合为空，商空间为零，取各 \(d_i=1\) 同样满足结论。
\end{proof}

例如 \(I=(x^2,xy,y^3)\subseteq k[x,y]\) 已由单项式生成；这组生成元是 Gröbner 基，因为理想中的每个单项式都被其中一个整除。标准单项式为 \(1,x,y,y^2\)，所以商的维数为四。只看生成元条数三，并不能得到这个维数。

\ProseParagraphBreak
维数只依赖商空间，标准单项式的具体选择却依赖单项式序。例如式\eqref{b1:eq:F-main-reduced-basis}给出的商，在 \(x\succ y\) 时使用 \(1,y\)，在 \(y\succ x\) 时使用 \(1,x\)。因为商中 \(x=y\)，两组坐标自然相容。更大规模的标准单项式及零维计算例子，可参阅 Sturmfels 的《Solving Systems of Polynomial Equations》\cite[Section 2.1]{Sturmfels2002PolynomialEquations}；乘法矩阵与重数的系统讨论留到第三卷第8.5节。

\subsection{商环的运算、单位与零因子}
\label{b1:subsec:F05:02}

令 \(\mathcal N\) 为标准单项式的线性生成空间。上一小节说明 \(\mathcal N\to R/I\) 是线性双射。由此在 \(\mathcal N\) 中表示商环运算：先按通常方式相加或相乘，再取正规余式。具体地，
\begin{equation}\label{b1:eq:F-normal-arithmetic}
\begin{aligned}
\operatorname{NF}_G(f+h)&=\operatorname{NF}_G(f)+\operatorname{NF}_G(h),\\
\operatorname{NF}_G(fh)&=
\operatorname{NF}_G\Bigl(\operatorname{NF}_G(f)\operatorname{NF}_G(h)\Bigr).
\end{aligned}
\end{equation}
第一式右端只含标准单项式，并与左端属于同一商类，所以相等。第二式中，两个内层余式分别与 \(f,h\) 模 \(I\) 同余，其乘积便与 \(fh\) 同余；再取唯一正规余式即可。特别地，第二式的最外层约化一般不能省去。

\begin{example}\label{b1:exm:F-two-point-algebra}
令
\[
I=(x-y,y^2-1),\qquad A=\mathbb Q[x,y]/I.
\]
在 \(A\) 中记 \(\bar y=y+I\)，则每个元素唯一写成 \(a+b\bar y\)。乘法为
\[
(a+b\bar y)(c+d\bar y)=(ac+bd)+(ad+bc)\bar y.
\]
取值同态
\begin{equation}\label{b1:eq:F-two-point-evaluation}
A\longrightarrow\mathbb Q\times\mathbb Q,
\qquad a+b\bar y\longmapsto(a+b,a-b)
\end{equation}
是代数同构，逆将 \((u,v)\) 送到 \((u+v)/2+(u-v)\bar y/2\)。两个幂等元
\[
e_+=\frac{1+\bar y}{2},\qquad e_-=\frac{1-\bar y}{2}
\]
满足 \(e_++e_-=1\)、\(e_+e_-=0\)。因此 \(a+b\bar y\) 可逆当且仅当 \(a+b\ne0\) 且 \(a-b\ne0\)；此时其逆为
\[
\frac{a-b\bar y}{a^2-b^2}.
\]
非零元素为零因子，当且仅当 \(a=b\) 或 \(a=-b\)。全部幂等元则为 \(0,e_+,e_-,1\)，因为 \(\mathbb Q\times\mathbb Q\) 的每个坐标只能是零或一。
\end{example}

这里计算的依据是正规式与明确的同构，而非由两个解的存在直接猜测商环。一般地，即使全部域值点都已求出，也可能遗漏商环中的非零幂零元素。

\begin{proposition}[单位判据]\label{b1:prop:F-unit-criterion}
设 \(k\) 为域，\(n\geq1\) 为整数，\(R=k[x_1,\ldots,x_n]\)，\(I\subsetneq R\) 为理想，\(f\in R\)。商类 \(f+I\) 在 \(R/I\) 中可逆，当且仅当 \(1\in I+(f)\)。若 \(I=(g_1,\ldots,g_s)\)，且存在 \(h,a_1,\ldots,a_s\in R\) 使
\[
1=hf+\sum_i a_i g_i,
\]
则其逆元为 \(h+I\)。在有效给定的系数域上，可以用 Gröbner 基的成员判定检验这个条件；计算过程中保留生成关系便能提取 \(h\)。
\end{proposition}
\begin{proof}
若 \((h+I)(f+I)=1+I\)，则 \(1-hf\in I\)，即 \(1\in I+(f)\)。反向，题示等式取商后给出 \((h+I)(f+I)=1+I\)，因环交换，这就是逆元。计算 Gröbner 基时，每个新增多项式都是此前多项式的组合；从初始生成元递推记录这些组合，最后把一的零余式表示代回，便得到所需等式。
\end{proof}

对前一例中的 \(2+\bar y\)，不必重跑完整算法，因为
\[
1=\frac{2-y}{3}(2+y)+\frac13(y^2-1).
\]
所以逆元是 \((2-\bar y)/3\)。这份等式可以直接展开核验，也是单位判据要求的具体证据。

\ProseParagraphBreak
再比较
\[
B=\mathbb Q[x,y]/(x,y^2).
\]
此处用 \(\bar y\) 表示 \(y+(x,y^2)\)。商环 \(B\) 同样以 \(1,\bar y\) 为基，维数为二，但 \(\bar y^2=0\)。元素 \(a+b\bar y\) 可逆当且仅当 \(a\ne0\)，逆为 \(a^{-1}-ba^{-2}\bar y\)。方程 \(x=0,y^2=0\) 在任何域扩张中只有一个点，而上例在特征零域中有两个点。这说明商环的向量空间维数不能一概当作不同解的个数；这里的区别可由非零幂零元直接看到，不必先建立一般重数理论。

\ProseParagraphBreak
本附录的单项式通过指数相加相乘，整除只需按各个指数比较。第二卷第20.6节的非交换字则必须保留字母次序；例如 \(xy\) 与 \(yx\) 不再自动相同，约化时要检查连续子字的包含与重叠。那里将用 Diamond 引理讨论唯一正规式。先掌握本附录的除法、首项和成员判定，再比较两种情形中哪些结论改变，会更容易理解非交换构造的条件。


\begin{chaptersummary}
单项式序给出了约化的下降方向，Gröbner 基则使原理想的全部首项都能由已知首项辨认。一般除法保证输出一个表示；对 Gröbner 基，输出的不可约余式才唯一，并且余式为零恰好等价于属于理想。有限生成性保证适合的基存在，Buchberger 算法进一步从给定生成组求出它。

\ProseParagraphBreak
S 多项式检测最高项抵消后可能出现的新关系。每个非零余式都保持原理想，却使首项理想严格增大；Hilbert 基定理给出的升链条件保证这种补充最终结束。约化基的唯一性要固定系数域和单项式序，不能把不同序下的输出直接逐项比较。

\ProseParagraphBreak
字典序让一部分关系不再含指定变量，标准单项式让每个商类有确定的坐标。消去所得关系的解未必全能提升，商空间维数也未必等于不同解的个数。保留生成关系、工作域和单项式序，才能把计算输出转成可以独立核验的结论。
\end{chaptersummary}

\BookExercises

\begin{exercise}\label{b1:ex:F-orders}
在 \(k[x,y,z]\) 中固定变量次序 \(x\succ y\succ z\)。分别用字典序、分次字典序、分次逆字典序排列 \(x^2z,xy^2,z^4\)，并求它们之和的首单项式。
\end{exercise}
\begin{solution}
字典序只先比较 \(x\) 指数，得到 \(x^2z\succ xy^2\succ z^4\)。分次字典序先比较全次数，得到 \(z^4\succ x^2z\succ xy^2\)。分次逆字典序同样把四次项放在最前；在两个三次项中，最后不同的 \(z\) 指数较小的项较大，故 \(z^4\succ xy^2\succ x^2z\)。三种首单项式分别为 \(x^2z,z^4,z^4\)。后两种次序的首项恰巧相同，但对剩下两项的比较不同。
\end{solution}

\begin{exercise}\label{b1:ex:F-nonwell-order}
在一元单项式上规定次数越小反而越大，即 \(1\succ x\succ x^2\succ\cdots\)。它与乘法相容吗？若按这个次序对一除以 \(1-x\)，逐次消去当前首项会发生什么？
\end{exercise}
\begin{solution}
两个指数比较后同时加同一个非负整数，大小关系保持，所以乘法相容。但所列就是无限严格下降链，故它不是良序。\(1-x\) 的首项是一；从一减去它后得到 \(x\)，再减去 \(x(1-x)\) 得 \(x^2\)，依次得到 \(x^m\)，永不终止。形式级数的无限求和是另一种运算，不能作为多项式除法的有限步输出。此例说明乘法相容性不能替代良序条件。
\end{solution}

\begin{exercise}\label{b1:ex:F-coefficient-field}
在 \(\mathbb Z[x]\) 中，\(x\) 的首单项式被 \(2x\) 的首单项式整除，为什么仍不能按本附录算法用 \(2x\) 消去 \(x\)？另在 \((\mathbb Z/6\mathbb Z)[x]\) 中计算 \((2x+1)(3x+1)\)，说明首单项式乘法公式的哪项基础可能失效。
\end{exercise}
\begin{solution}
消去首项需要系数 \(1/2\)，它不属于 \(\mathbb Z\)。而 \(x\notin(2x)\)，因为若 \(x=2xh\)，比较一次系数会要求某个整数的两倍等于一。扩域到 \(\mathbb Q\) 后才允许这个除法。

模六时乘积为 \(5x+1\)，首单项式为 \(x\)，不是两个首单项式乘积 \(x^2\)。最高项系数 \(2\cdot3=0\) 被零因子消去。故一般环既可能缺少首系数的逆，也可能缺少最高项乘积非零这一性质。
\end{solution}

\begin{exercise}\label{b1:ex:F-monomial-ideal}
对 \(J=(x^3y,x^2y^2,xy^3,x^4y^2,y^5)\subseteq k[x,y]\)，删去冗余单项式生成元，判断 \(x^2y^3,x^2y\) 是否属于 \(J\)，并判断商环是否有限维。
\end{exercise}
\begin{solution}
\(x^4y^2\) 被 \(x^3y\) 整除，可删去。余下四个单项式互不整除，所以不能再删。\(x^2y^3\) 被 \(x^2y^2\) 整除，属于 \(J\)；\(x^2y\) 不被任何生成元整除，不属于 \(J\)。所有生成元都含 \(y\)，故每个纯 \(x\) 幂都是标准单项式。由\cref{b1:thm:F-standard-monomial-basis}，商环无限维。
\end{solution}

\begin{exercise}\label{b1:ex:F-nonzero-remainder-member}
在 \(\mathbb Q[x,y]\) 的字典序 \(x\succ y\) 下，令 \(g_1=xy-1,g_2=y^2-2\)。证明 \(f=x^2-1/2\) 除以这个生成组时余式为自身，却仍属于 \((g_1,g_2)\)。
\end{exercise}
\begin{solution}
两个首单项式为 \(xy,y^2\)，都不整除 \(x^2\) 或一，所以除法没有可消去的项，余式就是 \(f\)。但是直接展开有
\[
x^2-\frac12=\frac{xy+1}{2}(xy-1)-\frac{x^2}{2}(y^2-2).
\]
故 \(f\) 属于原理想。非零余式能够否定成员资格的前提，是所用生成组已经为 Gröbner 基。
\end{solution}

\begin{exercise}\label{b1:ex:F-complete-second-example}
对上一题的 \(g_1,g_2\)，计算首个 S 多项式，求约化 Gröbner 基，并重新计算 \(x^2-1/2\) 的正规余式。
\end{exercise}
\begin{solution}
最小公倍式为 \(xy^2\)，所以 \(S(g_1,g_2)=yg_1-xg_2=2x-y\)。令 \(h=x-y/2\)。原理想等于 \((h,g_2)\)：一边由 \(2h=S(g_1,g_2)\) 得到，另一边由 \(g_1=yh+g_2/2\) 得到。\cref{b1:prop:F-triangular-basis}说明 \(\{x-y/2,y^2-2\}\) 是 Gröbner 基；两者首一，尾项不被任何首单项式整除，故已约化。代入 \(x=y/2\)，再模 \(y^2-2\)，得到 \(x^2-1/2\) 的余式为零，与前题证书相符。
\end{solution}

\begin{exercise}\label{b1:ex:F-reduced-versus-redundant}
对 \(I=(x-y,y^2)\subseteq\mathbb Q[x,y]\) 及字典序 \(x\succ y\)，判断 \(\{x-y,y^2,x^2\}\) 与 \(\{2x-2y,y^2\}\) 是否为 Gröbner 基、是否约化，并求 \(I\) 的约化基。
\end{exercise}
\begin{solution}
由\cref{b1:prop:F-triangular-basis}，\(G=\{x-y,y^2\}\) 为约化基。又 \(x^2=(x+y)(x-y)+y^2\in I\)，它的首单项式已属于 \((x,y^2)\)，所以把它加入 \(G\) 仍给出 Gröbner 基，却因首项冗余而不约化。第二组只是把第一项乘以非零常数，生成理想及首项理想均不变，也是 Gröbner 基；但第一项不首一，故不约化。两者约化后都得到 \(G\)。
\end{solution}

\begin{exercise}\label{b1:ex:F-coefficient-specialization}
分别在 \(\mathbb Q[x]\) 和 \(\mathbb F_2[x]\) 中计算由 \(x^2+1,x-1\) 生成的理想的约化 Gröbner 基。解释为何不能把有理数上的最终基逐项模二就当作第二个问题的答案。
\end{exercise}
\begin{solution}
在 \(\mathbb Q[x]\) 中，
\[
1=\tfrac12(x^2+1)-\tfrac12(x+1)(x-1),
\]
所以理想为全环，约化基为 \(\{1\}\)。在 \(\mathbb F_2[x]\) 中，\(x-1=x+1\)，且 \(x^2+1=(x+1)^2\)，所以理想为 \((x+1)\)，约化基为 \(\{x+1\}\)。前一个证书使用 \(1/2\)，没有可模二的意义。改变系数域时必须重新核对运算中所除的数。
\end{solution}

\begin{exercise}\label{b1:ex:F-change-order-cubic}
设 \(I=(x-y^2,y^3-1)\subseteq\mathbb Q[x,y]\)。求字典序 \(x\succ y\) 和 \(y\succ x\) 下的约化基，并分别给出商环的一组标准单项式。
\end{exercise}
\begin{solution}
第一种次序下由\cref{b1:prop:F-triangular-basis}，原两项已经是约化基，标准单项式为 \(1,y,y^2\)。商中 \(x=y^2,y^3=1\) 推出 \(x^2=y,x^3=1\)，故 \(y-x^2,x^3-1\) 属于 \(I\)。反过来，在按后两项取的商中，\(y=x^2,x^3=1\) 推出 \(y^2=x,y^3=1\)，所以两理想相同。交换三角生成组命题中的变量，便知第二种次序下的约化基为 \(\{y-x^2,x^3-1\}\)，标准单项式为 \(1,x,x^2\)。两组基都给出三维商空间，具体单项式却不同。
\end{solution}

\begin{exercise}\label{b1:ex:F-normal-degree-increase}
对 \(m\ge2\)，在 \(k[x,y]\) 中取字典序 \(x\succ y\) 及理想 \(I=(x-y^m)\)。证明 \(\{x-y^m\}\) 为 Gröbner 基，并求 \(x^d\)、\(d\ge1\)，的正规余式。这是否违反除法的终止性？
\end{exercise}
\begin{solution}
任一非零理想元素形如 \(h(x-y^m)\)，其首单项式为 \(\operatorname{LM}(h)x\)，所以被 \(x\) 整除，得到 Gröbner 基条件。反复用 \(x=y^m\) 得余式 \(y^{md}\)，其中没有可被 \(x\) 整除的项，故确为正规余式。它的全次数由 \(d\) 增大为 \(md\)。终止依据是指定单项式序的下降；在字典序中，任何 \(y\) 幂都小于 \(x\)，所以全次数增加并不违背这个依据。
\end{solution}


\begin{exercise}\label{b1:ex:F-coprime-leading}
设 \(f,g\) 首一，且首单项式 \(m,n\) 互素。证明 \(\{f,g\}\) 为 Gröbner 基。再对分次字典序 \(x\succ y\)，将结论用于 \(f=x^2+y,g=y^2+x\)，求商环的一组标准单项式。
\end{exercise}
\begin{solution}
写 \(f=m+f'\)、\(g=n+g'\)。因 \(\operatorname{lcm}(m,n)=mn\)，有
\[
S(f,g)=nf-mg=f'g-g'f.
\]
若 \(f'\ne0\)，则 \(\operatorname{LM}(f'g)=\operatorname{LM}(f')n\prec mn\)；另一项同理。这给出各项首单项式均严格低于最小公倍式的表示，正是\cref{b1:thm:F-buchberger-criterion}的\itemref{b1:item:F-buchberger-bound}，所以所给两项构成 Gröbner 基。尾项为零时对应求和项直接略去，结论仍成立。

应用中两个首单项式是 \(x^2,y^2\)，互素，所以已有 Gröbner 基。标准单项式为 \(1,x,y,xy\)，商环为四维。这里使用分次字典序；若把它换成单纯的字典序，\(y^2+x\) 的首单项式会改变，不能照搬此项判断。
\end{solution}

\begin{exercise}\label{b1:ex:F-elimination-order}
在 \(\mathbb Q[x,y]\) 中取 \(I=(x-y^2,y^2-2)\)。求 \(I\cap\mathbb Q[x]\) 和 \(I\cap\mathbb Q[y]\)。为什么不能从字典序 \(x\succ y\) 下的一组任意 Gröbner 基中，直接删去含 \(y\) 的项来求第一个交？
\end{exercise}
\begin{solution}
相加两个原生成元得到 \(x-2\)，反向有 \(x-y^2=(x-2)-(y^2-2)\)，所以 \(I=(x-2,y^2-2)\)。在两种变量优先次序下，这组多项式的首单项式分别来自不同变量，均为 Gröbner 基；也可直接用一元代入与除法证明两交为 \((x-2)\) 和 \((y^2-2)\)。

另一方面，\cref{b1:prop:F-triangular-basis}说明原两项本身也是 \(x\succ y\) 下的 Gröbner 基，但它们都含 \(y\)。若直接删去所有含 \(y\) 的多项式，会得到空集，错误地声称 \(I\cap\mathbb Q[x]=0\)。消去 \(y\) 的定理要求把 \(y\) 放在被优先比较的位置；所用次序的条件不能忽略。
\end{solution}

\begin{exercise}\label{b1:ex:F-cusp-parameter}
对任意域 \(k\)，在 \(k[t,x,y]\) 中令 \(J=(x-t^2,y-t^3)\)。证明 \(J\cap k[x,y]=(y^2-x^3)\)。若 \((a,b)\in K^2\)、\(K/k\) 为域扩张，满足 \(b^2=a^3\)，它是否总能写成 \((t^2,t^3)\)？
\end{exercise}
\begin{solution}
把 \(x,y\) 分别代为 \(t^2,t^3\) 的同态，其核为 \(J\)：对任一多项式先按 \(x-t^2\)、再按 \(y-t^3\) 整理，余下就是 \(k[t]\) 中的代入值。因此交理想是代入同态在 \(k[x,y]\) 上的核。

显然 \(y^2-x^3\) 在核中。把任意 \(f(x,y)\) 按关于 \(y\) 首一的 \(y^2-x^3\) 除，余式形如 \(A(x)+yB(x)\)。其代入值为 \(A(t^2)+t^3B(t^2)\)。前一项只含偶次幂，后一项只含奇次幂，单项式独立性使代入值为零当且仅当 \(A=B=0\)。故核没有更多元素。

对所给点，若 \(a=0\)，则 \(b=0\)，取 \(t=0\)。若 \(a\ne0\)，取 \(t=b/a\)，由 \(b^2=a^3\) 得 \(t^2=a\)、\(t^3=b\)。所以这里所有域值点都能提升，且上述公式对任意特征有效。
\end{solution}

\begin{exercise}\label{b1:ex:F-three-variable-projection}
在 \(k[x,y,z]\) 中令 \(I=(xy-1,yz-1)\)，用字典序 \(x\succ y\succ z\) 求一组 Gröbner 基及 \(I\cap k[z]\)。描述原方程的 \(k\) 值解在 \(z\) 坐标上的像。
\end{exercise}
\begin{solution}
令 \(f=xy-1,h=yz-1\)。计算
\[
zf-xh=x-z,
\qquad f=y(x-z)+h,
\]
所以 \(I=(x-z,yz-1)\)。这两项的 S 多项式为
\[
yz(x-z)-x(yz-1)=x-yz^2=(x-z)-z(yz-1),
\]
按所写基先消去 \(x\)，再消去 \(yz^2\)，余式为零。因此它们为 Gröbner 基。基中没有只含 \(z\) 的项，\cref{b1:thm:F-elimination}给出交理想为零。

原方程要求 \(z\ne0\)，此时唯一得到 \((x,y,z)=(z,z^{-1},z)\)。故投影像为 \(k^\times\)，而消去理想的 \(k\) 值零点是整个 \(k\)，仍多出一个不能提升的点。
\end{solution}

\begin{exercise}\label{b1:ex:F-two-equations}
在 \(\mathbb Q[x,y]\) 中取字典序 \(x\succ y\)，对方程组 \(x+y=1,x^2+y^2=1\) 求约化 Gröbner 基，给出全部有理解和复数解，并写明消去过程中所用的生成关系。
\end{exercise}
\begin{solution}
记 \(f=x+y-1\)、\(g=x^2+y^2-1\)。有
\[
g-(x-y+1)f=2(y^2-y).
\]
因二可逆，原理想等于 \((x+y-1,y^2-y)\)。由\cref{b1:prop:F-triangular-basis}，这已是 \(x\succ y\) 下的约化基。\(y^2-y=0\) 给出 \(y=0,1\)，再由 \(x=1-y\) 得 \((1,0),(0,1)\)。两点都在 \(\mathbb Q\) 中，且在 \(\mathbb C\) 中也没有更多根，所以两种解集相同。
\end{solution}

\begin{exercise}\label{b1:ex:F-five-dimensional-quotient}
对 \(A=k[x,y]/(x^3,x^2y,y^2)\)，求标准单项式基及维数。令 \(\mathfrak m=(\bar x,\bar y)\subseteq A\)，计算 \(\mathfrak m^2,\mathfrak m^3\)，并刻画全部单位。
\end{exercise}
\begin{solution}
标准单项式为 \(1,x,x^2,y,xy\)，所以维数为五。像理想 \(\mathfrak m\) 由后四项线性生成，其平方为 \(k\bar x^2+k\bar x\bar y\)，三次幂为零，因为所有三次单项式都被给定生成元整除。

每个元素唯一写为 \(a+u\)、\(a\in k,u\in\mathfrak m\)。若 \(a\ne0\)，逆为 \(a^{-1}-a^{-2}u+a^{-3}u^2\)，相乘后用 \(u^3=0\) 验证。若 \(a=0\)，则元素幂零，不能可逆。因此单位恰为常数项非零的元素。
\end{solution}

\begin{exercise}\label{b1:ex:F-two-normal-bases}
令 \(I=(x^2-y,y^2)\subseteq k[x,y]\)。分别用分次字典序 \(x\succ y\) 与字典序 \(y\succ x\) 求约化基和标准单项式，并证明所得商环同构于 \(k[t]/(t^4)\)。
\end{exercise}
\begin{solution}
分次字典序下，原首单项式为 \(x^2,y^2\)，其 S 多项式为 \(-y^3\)，可被 \(y^2\) 约化为零。因此原生成组为约化基，标准单项式是 \(1,x,y,xy\)。

字典序 \(y\succ x\) 下，理想也写成 \((y-x^2,x^4)\)：有 \(x^4=(x^2+y)(x^2-y)+y^2\)，反向 \(y^2=(y+x^2)(y-x^2)+x^4\)。交换\cref{b1:prop:F-triangular-basis}中的变量，得到 Gröbner 基 \(\{y-x^2,x^4\}\)；两项首一且尾项不可约，所以它已约化，标准单项式为 \(1,x,x^2,x^3\)。赋值 \(t\mapsto\bar x\)，以及反向 \(\bar x\mapsto t,\bar y\mapsto t^2\)，定义互逆同态，故得到所述同构。两组基的对应是 \(\bar y=\bar x^2,\bar x\bar y=\bar x^3\)。
\end{solution}

\begin{exercise}\label{b1:ex:F-inverse-certificate}
在 \(A=\mathbb Q[x,y]/(x-y^2,y^3-2)\) 中，求 \(x+y+1\) 的剩余类的逆，并将结果写成原多项式环内的一条单位证书。
\end{exercise}
\begin{solution}
商中该元素等于 \(y^2+y+1\)，而
\((y-1)(y^2+y+1)=y^3-1\equiv1\)。故逆为 \(y-1\) 的剩余类。完整证书是
\[
1=(y-1)(x+y+1)-(y-1)(x-y^2)-(y^3-2).
\]
展开后立即得到一；由\cref{b1:prop:F-unit-criterion}，它同时验证原元素确为单位及所给逆元正确。
\end{solution}

\begin{exercise}\label{b1:ex:F-characteristic-two-main}
把本附录的双点例子改到 \(\mathbb F_2\)：求 \(A=\mathbb F_2[x,y]/(x-y,y^2-1)\) 的维数、全部单位和幂等元。它在任意域扩张中有几个不同的解？
\end{exercise}
\begin{solution}
令 \(z=\bar y-1\)，则 \(z^2=\bar y^2-2\bar y+1=0\)，且 \(\bar x=\bar y=1+z\)。因此 \(A\cong\mathbb F_2[z]/(z^2)\)，维数为二，四个元素是 \(0,1,z,1+z\)。单位为 \(1,1+z\)，因为常数项必须非零，并且 \((1+z)^2=1\)。若 \(a+bz\) 幂等，则其平方为 \(a\)，所以 \(b=0\)，只剩零与一。原方程的根满足 \((y-1)^2=0\)，在任何域中都迫使 \(y=x=1\)，故只有一个不同的解。有理数情形中除以二得到的两个幂等元，在这里已不能照写。
\end{solution}

\begin{exercise}\label{b1:ex:F-points-not-membership}
给出下列两种现象的直接证明：在 \(k[x]\) 中，\(x\) 在理想 \((x^2)\) 的全部域值零点上为零，却不属于该理想；在 \(\mathbb F_2[x]\) 中，非零多项式 \(x^2-x\) 在每个 \(\mathbb F_2\) 元素上都为零。它们为何不能代替 Gröbner 理想成员判定？
\end{exercise}
\begin{solution}
在任意域扩张中，\(a^2=0\) 蕴含 \(a=0\)，所以第一项确在全部零点消失。但 \(x\) 除以 \(x^2\) 的余式为 \(x\ne0\)，故不属于 \((x^2)\)。第二项代入零、一都为零，却具有非零的多项式系数，因而不属于零理想。前一例即使允许全部域扩张的点，仍没有检测到幂零关系的次数；后一例只考察有限基域中的点，连多项式是否为零都不能判定。成员判据检查的是多项式关系，而不是有限个取值。
\end{solution}

\begin{exercise}\label{b1:ex:F-truncated-total-degree}
设 \(d\ge1\)，在 \(k[x,y]\) 中令 \(I=(x,y)^d\)。求一个对任意全局单项式序都适用的 Gröbner 基，并求商环维数、其中像理想 \((\bar x,\bar y)\) 的幂零指数及全部单位。
\end{exercise}
\begin{solution}
理想由全次数恰为 \(d\) 的单项式 \(x^d,x^{d-1}y,\ldots,y^d\) 生成。这些单项式构成任意单项式序下的 Gröbner 基。标准单项式恰为全次数小于 \(d\) 的单项式，因此维数为
\[
\sum_{j=0}^{d-1}(j+1)=\frac{d(d+1)}2.
\]
像理想的 \(d\) 次幂为零；若 \(d>1\)，其 \(d-1\) 次幂含非零标准单项式 \(\bar x^{d-1}\)，所以幂零指数恰为 \(d\)。\(d=1\) 时像理想本身为零，最小正指数也为一。写元素为 \(a+u\)，其中 \(u\) 在像理想中；\(a\ne0\) 时逆为有限和 \(a^{-1}\sum_{j=0}^{d-1}(-a^{-1}u)^j\)，\(a=0\) 时元素幂零，不能可逆。
\end{solution}

\begin{exercise}\label{b1:ex:F-ideal-comparisons}
在 \(\mathbb Q[x,y]\) 中比较
\[
I=(x-y,y^2-1),\quad J=(x-y,x^2-1),\quad K=(x-y,y^3-y).
\]
判断哪些相等、哪些包含，并用正规余式给出严格包含的见证。
\end{exercise}
\begin{solution}
由 \(x^2-1=(x+y)(x-y)+(y^2-1)\)，以及该式移项后的反向表达，得到 \(I=J\)。又 \(y^3-y=y(y^2-1)\)，故 \(K\subseteq I\)。\cref{b1:prop:F-triangular-basis}说明 \(\{x-y,y^3-y\}\) 为 \(K\) 的 Gröbner 基；\(y^2-1\) 对它的正规余式仍为自身，非零，因此 \(y^2-1\in I\setminus K\)。所以 \(K\subsetneq I=J\)。只知一组生成元在另一理想中还不够判断相等，必须检查反向包含。
\end{solution}
